\documentclass[12pt,a4paper,UTF8]{ctexart}
\usepackage[a4paper,top=22mm,bottom=22mm,left=22mm,right=22mm,headheight=15pt]{geometry}
\newcommand{\Needspace}[1]{\par\begingroup\skip0=\lastskip\vskip-\skip0\dimen0=#1\relax\advance\dimen0 by\skip0\vskip\dimen0\penalty-100\vskip-\dimen0\vskip\skip0\penalty9999\endgroup}
\usepackage{amsmath,amssymb,mathtools,bm}
\usepackage{graphicx,xcolor,booktabs,longtable,array,enumitem,fancyhdr,caption,float,microtype,setspace}
\usepackage[hidelinks,unicode]{hyperref}
\usepackage[most]{tcolorbox}
\definecolor{ContestBlue}{HTML}{17365D}
\hypersetup{pdftitle={第079套},pdfauthor={物理竞赛教学资料}}
\setstretch{1.24}
\setlength{\parindent}{2em}
\setlength{\parskip}{0.35em}
\setlength{\emergencystretch}{2em}
\setlist{itemsep=0.35em,topsep=0.4em}
\captionsetup{font=small,labelfont=bf,justification=centering}
\setcounter{secnumdepth}{0}
\pagestyle{fancy}
\fancyhf{}
\fancyhead[L]{\small 第079套}
\fancyfoot[C]{\thepage}
\renewcommand{\headrulewidth}{0.4pt}
\ctexset{
  section={format=\Large\bfseries\color{ContestBlue},beforeskip=1.2em,afterskip=0.7em},
  subsection={format=\large\bfseries\color{ContestBlue},beforeskip=1.0em,afterskip=0.5em},
  subsubsection={format=\normalsize\bfseries,beforeskip=0.8em,afterskip=0.35em}
}
\usepackage{calc}
\newcounter{none}
\let\originalsubsection\subsection
\renewcommand{\subsection}{\Needspace{6\baselineskip}\originalsubsection}
\makeatletter
\newcommand{\tightlist}{\setlength{\itemsep}{0pt}\setlength{\parskip}{0pt}}
\newcommand{\maxwidth}{\ifdim\Gin@nat@width>\linewidth\linewidth\else\Gin@nat@width\fi}
\newcommand{\maxheight}{\ifdim\Gin@nat@height>0.72\textheight 0.72\textheight\else\Gin@nat@height\fi}
\newcommand{\pandocbounded}[1]{\begingroup\setkeys{Gin}{width=\maxwidth,height=\maxheight,keepaspectratio}#1\endgroup}
\makeatother
\begin{document}
\begin{titlepage}
  \centering
  \vspace*{58mm}
  {\Huge\bfseries\color{ContestBlue} 第079套\par}
  \vfill
  \vspace*{28mm}
\end{titlepage}
\section{第 1 题}\label{ux7b2c-1-ux9898}

两根完全相同的四分之一细圆弧由均匀金属制成，具有弹性。两根圆弧分别按图示方式将一端固定在墙上，另一端各悬挂一个质量相同的重物。形变很小，圆弧质量可以忽略。比较两种固定方式中重物下降的高度。

\begin{figure}[H]
\centering
\includegraphics[width=0.85\linewidth,height=\textheight,keepaspectratio,alt={图 1-1　两种四分之一圆弧的固定方式及相同的竖直载荷。}]{../figures/question-01.pdf}
\caption*{图 1-1　两种四分之一圆弧的固定方式及相同的竖直载荷。}
\end{figure}

\subsection{一、题型定位}\label{ux4e00ux9898ux578bux5b9aux4f4d}

两种安装的固定端弯矩相同，沿弧的力臂分布却不同。以同一角坐标逐截面比较，再用弯曲能或单位载荷法求自由端下降量。

\subsection{二、知识点}\label{ux4e8cux77e5ux8bc6ux70b9}

\begin{itemize}
\tightlist
\item
  平面曲梁截面的弯矩分布。
\item
  小变形下的线弹性关系。
\item
  弯曲应变能 \(U=\int M^2/(2EI)\,ds\)。
\item
  卡氏定理与载荷方向位移。
\item
  无量纲积分的大小比较。
\end{itemize}

\subsection{三、物理图像}\label{ux4e09ux7269ux7406ux56feux50cf}

重物的拉力均为
\(P=mg\)。对任一截面，决定弯曲程度的是重力作用线到该截面的垂直距离，而不是截面到自由端的弧长。两种安装方式的力臂函数不同，因此即使圆弧长度、材料和端点载荷相同，储存的弯曲能也不同。

\subsection{四、解题入口}\label{ux56dbux89e3ux9898ux5165ux53e3}

令圆弧半径为 \(R\)，弯曲刚度为 \(EI\)。从自由端起沿圆弧取角坐标
\(\varphi\in[0,\pi/2]\)。分别写出两种方式的截面弯矩，再利用

\[U=\int_0^{\pi R/2}\frac{M^2}{2EI}\,ds,
\qquad
h=\frac{\partial U}{\partial P}.\]

\subsection{五、解析}\label{ux4e94ux89e3ux6790}

\begin{figure}[H]
\centering
\includegraphics[width=0.9\linewidth,height=\textheight,keepaspectratio,alt={图 1-2　两种安装方式的截面力臂和弯矩函数。}]{../figures/solution-01.pdf}
\caption*{图 1-2　两种安装方式的截面力臂和弯矩函数。}
\end{figure}

\Needspace{3\baselineskip}
\textbf{1. 方案 (a) 的弯矩}

图 1-2 左侧中，截面到重力作用线的水平距离为 \(R\sin\varphi\)，所以

\[M_a(\varphi)=PR\sin\varphi.\]

圆弧微元长度为

\[ds=R\,d\varphi.\]

故方案 (a) 的弯曲应变能为

\[
\begin{aligned}
U_a
&=\int_0^{\pi/2}\frac{M_a^2}{2EI}R\,d\varphi\\
&=\frac{P^2R^3}{2EI}\int_0^{\pi/2}\sin^2\varphi\,d\varphi.
\end{aligned}
\]

由

\[\sin^2\varphi=\frac{1-\cos2\varphi}{2},\]

得到

\[
\int_0^{\pi/2}\sin^2\varphi\,d\varphi
=\left[\frac{\varphi}{2}-\frac{\sin2\varphi}{4}\right]_0^{\pi/2}
=\frac{\pi}{4}.
\]

因此

\[U_a=\frac{\pi P^2R^3}{8EI}.\]

\Needspace{3\baselineskip}
\textbf{2. 方案 (b) 的弯矩}

图 1-2 右侧中，截面到重力作用线的水平距离为 \(R(1-\cos\varphi)\)，所以

\[M_b(\varphi)=PR(1-\cos\varphi).\]

相应的弯曲应变能为

\[
U_b=\frac{P^2R^3}{2EI}
\int_0^{\pi/2}(1-\cos\varphi)^2\,d\varphi.
\]

把平方展开：

\[
(1-\cos\varphi)^2
=1-2\cos\varphi+\cos^2\varphi.
\]

逐项积分得到

\[
\begin{aligned}
\int_0^{\pi/2}(1-\cos\varphi)^2\,d\varphi
&=\frac{\pi}{2}-2+\frac{\pi}{4}\\
&=\frac{3\pi}{4}-2.
\end{aligned}
\]

所以

\[U_b=\frac{P^2R^3}{2EI}\left(\frac{3\pi}{4}-2\right).\]

\Needspace{3\baselineskip}
\textbf{3. 由卡氏定理求竖直下降量}

因为 \(U_a\)、\(U_b\) 都与 \(P^2\) 成正比，载荷作用点沿载荷方向的位移为

\[h_i=\frac{\partial U_i}{\partial P},\qquad i=a,b.\]

于是

\[h_a=\frac{\pi PR^3}{4EI},\]

\[h_b=\frac{PR^3}{EI}\left(\frac{3\pi}{4}-2\right).\]

比较两个无量纲系数：

\[
\frac{\pi}{4}-\left(\frac{3\pi}{4}-2\right)
=2-\frac{\pi}{2}>0.
\]

因此

\[\boxed{h_a>h_b}.\]

也就是说，图 1-1 中方案 (a) 的自由端下降得更多。

\Needspace{3\baselineskip}
\textbf{逐截面比较的另一种写法}

在 \(0\le\varphi\le\pi/2\) 内，

\[
\sin\varphi-(1-\cos\varphi)
=2\sin\frac\varphi2\left(\cos\frac\varphi2-\sin\frac\varphi2\right)\ge0.
\]

因此
\(M_a(\varphi)\ge M_b(\varphi)\)，且除端点外严格大于。线弹性条件下曲率改变与弯矩成正比，方案
(a) 在绝大多数截面弯得更多。这个局部判断与能量积分所得结论一致。

\Needspace{3\baselineskip}
\textbf{为什么不能只比较固定端弯矩}

两种方式在固定端的最大弯矩都等于
\(PR\)。如果只看固定端，会误以为两种形变相同。自由端位移取决于整根圆弧上
\(M^2\) 的积分；固定端只是一个截面，其信息不足以决定整体柔顺度。

\Needspace{3\baselineskip}
\textbf{角坐标和薄杆近似的核对}

图 1-2 中 \(\varphi\) 都从自由端开始：方案 (a) 用竖直半径为零角，方案
(b) 用向左水平半径为零角。两图中的水平投影分别为 \(R\sin\varphi\) 与
\(R(1-\cos\varphi)\)。自由端 \(\varphi=0\) 时两弯矩均为零；固定端
\(\varphi=\pi/2\) 时均为 \(PR\)，这两个端点可先验算。

细杆模型忽略伸缩、剪切柔度。若计入轴力和剪力，两种姿态的
\(\sin^2\varphi\)、\(\cos^2\varphi\) 积分均为
\(\pi/4\)，因此轴向与剪切对竖向柔度的附加贡献相同，仍有
\(h_a>h_b\)；但总位移比不再恰为
\(2.20\)。上文的比值只指以弯曲为主的细梁结果。

任意角 \(\beta\) 的两积分在 \(0<\beta<\pi/2\)
都可逐点比较；\(\beta\to0\) 时
\(h_a/h_b\sim20/(3\beta^2)\)。水平载荷变式须先重新画水平力的竖直力臂，不能只把原公式里的
\(P\) 改个方向。

\Needspace{3\baselineskip}
\textbf{伸缩柔度的共同项可以直接积分}

若截面面积为 \(A_c\)，两种安装的轴力分别按角坐标含
\(P\cos\varphi\)、\(P\sin\varphi\)。轴向能的载荷导数均为

\[h_{\rm axial}=\frac{PR}{EA_c}\int_0^{\pi/2}\sin^2\varphi,d\varphi
=\frac{\pi PR}{4EA_c}.\]

剪切修正采用同一截面的剪切柔度系数 \(1/(k_sGA_c)\)，也给共同项
\(\pi PR/(4k_sGA_c)\)。所以总位移可写成

\[h_a=P(C_{a,\rm bend}+C_0),\qquad h_b=P(C_{b,\rm bend}+C_0),\]
\[\frac{h_a}{h_b}=\frac{C_{a,\rm bend}+C_0}{C_{b,\rm bend}+C_0}.\]

共同项 \(C_0>0\)
会使比值从2.205向1靠近，但位移差仍为两弯曲柔度之差。若误把伸缩项只加给某一个方案，便可能错误反转刚度顺序。这里仍沿用薄梁、小变形、线弹性；厚曲梁应重新确定应力分布，不能仅添加两项就称为精确厚梁解。

\Needspace{3\baselineskip}
\textbf{用单位载荷法独立求位移}

在重物处再施加竖直单位力，各截面的单位弯矩分别为
\(m_a=R\sin\varphi\)、\(m_b=R(1-\cos\varphi)\)。单位载荷法给

\[h_i=\int\frac{M_i m_i}{EI}ds.\]

由于真实载荷是 \(P\)，\(M_i=Pm_i\)，所以

\[h_a=\frac{PR^3}{EI}\int_0^{\pi/2}\sin^2\varphi\,d\varphi
=\frac{\pi PR^3}{4EI},\]

\[h_b=\frac{PR^3}{EI}\int_0^{\pi/2}(1-\cos\varphi)^2d\varphi
=\frac{(3\pi-8)PR^3}{4EI}.\]

这与能量求导一致，却不需要先记住系数 \(1/2\)。比值
\(\pi/(3\pi-8)\simeq2.205\)；两者差值为
\((2-\pi/2)PR^3/(EI)>0\)。自由端弯矩为零，固定端都为
\(PR\)，但中间弯矩分布不同，所以只比较固定端最大弯矩不能判断位移大小。

以上 \(P\) 是力而不是质量，若题给重物质量 \(m\)，最后才代
\(P=mg\)。\(EI\) 的量纲为力乘长度平方，\(PR^3/(EI)\)
为长度，符合下降量的单位。

\Needspace{3\baselineskip}
\textbf{小圆弧极限可检验哪个方案更硬}

把圆心角改为 \(\beta\)，保持同样的安装和角坐标。两个柔度积分为

\[J_a(\beta)=\frac\beta2-\frac{\sin2\beta}{4},\]
\[J_b(\beta)=\frac{3\beta}{2}-2\sin\beta+\frac{\sin2\beta}{4}.\]

小角展开得到 \(J_a\sim\beta^3/3\)、\(J_b\sim\beta^5/20\)，故
\(h_a/h_b\sim20/(3\beta^2)\)。方案b的载荷线在短弧内更接近各截面，所以弯曲力臂是二阶小量；方案a的力臂是一阶小量。这一几何解释与两个展开的阶数对应。

若把重物沿另一个方向拉，单位载荷也必须沿新的测量方向；想求某截面的转角，则引入单位力矩，不能继续使用竖直单位力的
\(m_i\)。若圆弧很厚或伸缩不能忽略，弯曲积分外还要加 \(\int Nn/(EA)ds\)
和剪切项，2.205便只是弯曲部分的比值。

\newpage

\section{第 2 题}\label{ux7b2c-2-ux9898}

一圆柱形电容器外半径为 \(R\)，内半径为 \(R-d\)，其中 \(d\ll R\)；长度为
\(l\)，总质量为
\(M\)。电容器由三根对称分布的绝缘细线悬挂，只能绕中心轴转动而不能水平移动。空间中有竖直向下、磁感应强度为
\(B\) 的匀强磁场。电容器已经充电，两极板间电势差为 \(V\)。

分别分析以下过程产生的转动：

\begin{enumerate}
\def\labelenumi{\arabic{enumi}.}
\tightlist
\item
  保持磁场不变，用位于端面的径向导线组将两极板连接，使电容器放电；
\item
  保持电容器带电，令磁场突然减小到零。
\end{enumerate}

\begin{figure}[H]
\centering
\includegraphics[width=0.846\linewidth,height=\textheight,keepaspectratio,alt={图 2-1　轴向磁场中的薄间隙圆柱电容器及其悬挂方式。}]{../figures/question-02.pdf}
\caption*{图 2-1　轴向磁场中的薄间隙圆柱电容器及其悬挂方式。}
\end{figure}

\subsection{一、题型定位}\label{ux4e00ux9898ux578bux5b9aux4f4d-1}

把两过程的力矩都换成角冲量，便可绕开未知放电速度和磁场变化时间。最后统一正方向，比较它们传给圆筒的角动量。

\subsection{二、知识点}\label{ux4e8cux77e5ux8bc6ux70b9-1}

\begin{itemize}
\tightlist
\item
  薄间隙同轴电容器的电容。
\item
  洛伦兹力与角冲量。
\item
  法拉第电磁感应定律。
\item
  电磁场角动量与机械角动量转化。
\item
  薄壁圆柱绕轴转动惯量。
\end{itemize}

\subsection{三、物理图像}\label{ux4e09ux7269ux7406ux56feux50cf-1}

径向运动的电荷在轴向磁场中受到切向洛伦兹力，因此放电过程会给导体一个轴向角冲量。若电荷不动而磁场变化，空间中出现环向电场；两极板上的异号电荷受到方向配合的切向力，同样获得轴向角动量。

\subsection{四、解题入口}\label{ux56dbux89e3ux9898ux5165ux53e3-1}

先求电容和电荷量。薄间隙近似下

\[C\simeq\frac{2\pi\varepsilon_0Rl}{d},
\qquad Q=CV.\]

然后分别计算两种过程的总角冲量，最后除以绕轴转动惯量 \(I\simeq MR^2\)。

\subsection{五、解析}\label{ux4e94ux89e3ux6790-1}

\begin{figure}[H]
\centering
\includegraphics[width=0.9\linewidth,height=\textheight,keepaspectratio,alt={图 2-2　径向电荷运动在轴向磁场中产生切向冲量。}]{../figures/solution-02.pdf}
\caption*{图 2-2　径向电荷运动在轴向磁场中产生切向冲量。}
\end{figure}

\Needspace{3\baselineskip}
\textbf{1. 电容器所带电荷}

同轴电容器的精确电容为

\[C=\frac{2\pi\varepsilon_0l}{\ln\dfrac{R}{R-d}}.\]

因为 \(d\ll R\)，

\[
\ln\frac{R}{R-d}
=-\ln\left(1-\frac dR\right)
\simeq\frac dR.
\]

故

\[
C\simeq\frac{2\pi\varepsilon_0Rl}{d},
\]

\[
Q\simeq\frac{2\pi\varepsilon_0RlV}{d}.
\]

\Needspace{3\baselineskip}
\textbf{2. 径向放电产生的角冲量}

设一小份正电荷 \(dq\) 从半径 \(R\) 运动到 \(R-d\)。其径向速度为
\(v_r\)，所受切向磁力大小为

\[dF_\varphi=dq\,v_rB.\]

对轴的力矩为

\[d\tau_z=r\,dF_\varphi=dq\,B,r v_r.\]

利用 \(v_r=dr/dt\)，这一小份电荷获得的角冲量为

\[
\begin{aligned}
dL_z
&=\int d\tau_z\,dt\\
&=dq\,B\int_R^{R-d}r\,dr.
\end{aligned}
\]

取大小：

\[
|dL_z|
=\frac{dqB}{2}\left[R^2-(R-d)^2\right].
\]

将平方差展开：

\[R^2-(R-d)^2=2Rd-d^2\simeq2Rd.\]

所以

\[|dL_z|\simeq dq\,BRd.\]

对全部电荷积分得到

\[|L_z|\simeq QBRd.\]

磁力对载流子的切向作用通过导线传给圆筒，圆筒获得相同方向的机械角冲量，大小为
\(QBRd\)。这里不能把载流子与导线说成获得两个相反角动量；磁场源与场也参与总角动量账目。以
\(I\simeq MR^2\) 近似，有

\[
\omega_1=\frac{QBRd}{MR^2}.
\]

代入 \(Q\)：

\[
\begin{aligned}
\omega_1
&=\frac{1}{MR^2}
\frac{2\pi\varepsilon_0RlV}{d}BRd\\
&=\boxed{\frac{2\pi\varepsilon_0BlV}{M}}.
\end{aligned}
\]

\Needspace{3\baselineskip}
\textbf{3. 磁场突然消失时的涡旋电场}

对半径为 \(r\) 的同心圆回路应用法拉第定律：

\[
\oint\bm E\cdot d\bm l
=2\pi rE_\varphi(r)
=-\frac{d}{dt}(\pi r^2B).
\]

因此

\[E_\varphi(r)=-\frac r2\frac{dB}{dt}.\]

外极板带 \(+Q\)，内极板带 \(-Q\)。两极板所受总力矩为

\[
\tau_z
=QRE_\varphi(R)-Q(R-d)E_\varphi(R-d).
\]

代入环向电场：

\[
\begin{aligned}
\tau_z
&=-\frac Q2\frac{dB}{dt}
\left[R^2-(R-d)^2\right]\\
&\simeq-QRd\frac{dB}{dt}.
\end{aligned}
\]

从初始磁场 \(B\) 积分到零：

\[
\Delta L_z
=\int\tau_zdt
=-QRd\int_B^0dB
=QBRd.
\]

故

\[
\boxed{\omega_2=\frac{QBRd}{MR^2}
=\frac{2\pi\varepsilon_0BlV}{M}}.
\]

两种过程中角速度大小相同；实际转向由极板正负与 \(\bm B\)
方向用右手定则确定。

\Needspace{3\baselineskip}
\textbf{场角动量视角}

充电同轴电容器的径向电场与轴向磁场共存时，坡印廷矢量
\(\bm S=(1/\mu_0)\bm E\times\bm B\)
沿环向，电磁场携带轴向角动量。无论让电荷消失还是让磁场消失，初始场角动量都必须转移给物质，所以两种最终机械角动量相同。

\Needspace{3\baselineskip}
\textbf{近似的边界}

\(d\ll R\)
同时用于三处：电容近似、平方差近似和转动惯量近似。若质量在内、外筒间分布不均，必须用实际
\(I\) 替换 \(MR^2\)，但角冲量 \(QBRd\) 的一阶结果不变。

\Needspace{3\baselineskip}
\textbf{带符号的一套表达式}

取从上方看逆时针为 \(+z\)，外筒 \(+Q\)，轴向磁场为 \(B_z\)；题图向下对应
\(B_z=-B\)。正电荷沿导线由外向内，磁力环向分量为
\(F_\varphi=-qB_z\dot r\)。因此

\[\Delta L_z=-QB_z\int_R^{R-d}r\,dr
=\frac{QB_z}{2}[R^2-(R-d)^2].\]

磁场消失时 \(E_\varphi=-r\dot B_z/2\)，

\[\Delta L_z=-\frac Q2[R^2-(R-d)^2]\int_{B_z}^{0}dB_z
=\frac{QB_z}{2}[R^2-(R-d)^2].\]

两者均负，即从上方看顺时针。解图分别标外筒、内筒相反感应力矩，净方向由外筒决定。

若质量分别集中在外、内薄筒，\(I_z=M_oR^2+M_i(R-d)^2\)。令
\(f=M_i/M\)、\(\delta=d/R\)，连同精确电容展开，

\[\omega=\frac{2\pi\varepsilon_0BlV}{M}[1+(2f-1)\delta+O(\delta^2)].\]

两筒质量相等时一次修正抵消；只保留角冲量的
\(d^2\)、却漏电容或惯量的同阶项，会算错修正。

\Needspace{3\baselineskip}
\textbf{厚间隙例子检验被消去的几何参数}

只作核算例，取内半径 \(b=R/2\)，内外筒质量各
\(M/2\)。这已经不属薄间隙。此时

\[I_z=\frac58MR^2,\qquad
C=\frac{2\pi\varepsilon_0l}{\ln2},\qquad
|\Delta L_z|=\frac38QBR^2.\]

故

\[|\omega|=\frac65\frac{\pi\varepsilon_0BlV}{M\ln2}.\]

与薄间隙值 \(2\pi\varepsilon_0BlV/M\) 的比为
\(3/(5\ln2)\simeq0.866\)。两筒等质量虽使薄间隙的一次修正抵消，并不让所有高阶修正也消失。这个有限几何例子能检查``与
\(d\) 无关''的适用范围。

保持电压而不是保持电荷会使 \(Q\) 随间隙改变；若题改为给定
\(Q\)，应从角冲量式直接代入，不能沿用包含精确电容的固定电压结论。两种实验控制条件产生不同的尺寸依赖。

\Needspace{3\baselineskip}
\textbf{不作薄间隙近似，两个过程仍相同}

写 \(b=R-d\)。同轴电容器的精确电容和电荷为

\[C=\frac{2\pi\varepsilon_0l}{\ln(R/b)},\qquad Q=CV.\]

两过程的角冲量绝对值都是 \(QB(R^2-b^2)/2\)，因此

\[|\omega|=\frac{\pi\varepsilon_0BlV(R^2-b^2)}{I_z\ln(R/b)}.\]

利用
\(R^2-b^2\simeq2Rd\)、\(\ln(R/b)\simeq d/R\)、\(I_z\simeq MR^2\)，才得到
\(2\pi\varepsilon_0BlV/M\)。薄间隙公式里 \(R,d\) 消失，是电容的 \(R/d\)
与角冲量的 \(Rd\) 相消，并不表示任意厚度下答案与半径和质量分布无关。

改变放电电阻会改变电流随时间的分布，却不改变流过的总电量
\(Q\)；只要磁场恒定、初终状态相同，角冲量不变。若放电和撤场同时发生，电荷量与磁场都随时间变化，应将磁力矩和感应电力矩同时积分，不能把两个完成过程的角冲量直接相加。

\Needspace{3\baselineskip}
\textbf{一次修正需同时保留三处展开}

令 \(\delta=d/R\)、\(f=M_i/M\)。分别有

\[\ln\frac1{1-\delta}=\delta+\frac{\delta^2}{2}+O(\delta^3),\]
\[R^2-b^2=2R^2\delta(1-\delta/2),\]
\[I_z=MR^2[1-2f\delta+O(\delta^2)].\]

把三者代入精确式，

\[|\omega|=\frac{2\pi\varepsilon_0BlV}{M}
[1+(2f-1)\delta+O(\delta^2)].\]

若内外筒质量相等，\(f=1/2\)，一次项为零；若质量主要在外筒，一次项为负。这个差别无法由薄壁总惯量
\(MR^2\)
看出。撤场时两筒受相反力矩，外筒的力矩绝对值较大；图中的两个切向箭头应分属内外筒，再对同一轴求和。

\newpage

\section{第 3 题}\label{ux7b2c-3-ux9898}

质量为 \(m\)、半径为 \(R\)
的乒乓球斜向下撞击水平桌面。碰撞前速度的水平分量为 \(v_0>0\)，角速度为
\(\omega_0\)；转轴水平且垂直于速度的水平分量。规定上旋
\(\omega>0\)、下旋 \(\omega<0\)，反弹后水平速度与 \(v_0\) 同向为正。

作如下假设：竖直方向碰撞完全弹性；接触期间的滑动距离可忽略；离开桌面瞬间接触点无滑动；忽略空气阻力。已知球绕相应直径的转动惯量为

\[I=\frac23mR^2.\]

分别讨论 \(\omega_0=0\)、\(\omega_0>0\)、\(\omega_0<0\) 时反弹瞬间的
\(v_f\)、\(\omega_f\) 及符号条件。

\begin{figure}[H]
\centering
\includegraphics[width=0.9\linewidth,height=\textheight,keepaspectratio,alt={图 3-1　乒乓球碰撞前后的水平速度和自转方向。}]{../figures/question-03.pdf}
\caption*{图 3-1　乒乓球碰撞前后的水平速度和自转方向。}
\end{figure}

\subsection{一、题型定位}\label{ux4e00ux9898ux578bux5b9aux4f4d-2}

一个切向冲量同时改变水平速度与自转。先用离桌瞬间的无滑动条件闭合方程，再判断弱上旋、强上旋和下旋的反弹方向。

\subsection{二、知识点}\label{ux4e8cux77e5ux8bc6ux70b9-2}

\begin{itemize}
\tightlist
\item
  线冲量---动量定理。
\item
  角冲量---角动量定理。
\item
  接触点相对速度。
\item
  滚动约束 \(v_f=R\omega_f\)。
\item
  分段符号判断与临界条件。
\end{itemize}

\subsection{三、物理图像}\label{ux4e09ux7269ux7406ux56feux50cf-2}

碰撞极短，桌面对球的切向作用可用冲量 \(J\)
表示。这个冲量既改变质心水平动量，也通过力臂 \(R\)
改变自转。题设的末态无滑动不是全过程纯滚动，而只是离桌瞬间的闭合条件。

\subsection{四、解题入口}\label{ux56dbux89e3ux9898ux5165ux53e3-2}

统一规定向右的切向冲量为正。列三式：

\[m(v_f-v_0)=J,\]

\[I(\omega_f-\omega_0)=-JR,\]

\[v_f=R\omega_f.\]

\subsection{五、解析}\label{ux4e94ux89e3ux6790-2}

\begin{figure}[H]
\centering
\includegraphics[width=0.732\linewidth,height=\textheight,keepaspectratio,alt={图 3-2　切向冲量对质心平动和绕心转动的同步改变。}]{../figures/solution-03.pdf}
\caption*{图 3-2　切向冲量对质心平动和绕心转动的同步改变。}
\end{figure}

\Needspace{3\baselineskip}
\textbf{1. 消去碰撞冲量}

由平动冲量方程

\[J=m(v_f-v_0).\]

代入转动方程：

\[
I(\omega_f-\omega_0)
=-mR(v_f-v_0).
\]

再代入 \(I=\frac23mR^2\)：

\[
\frac23mR^2(\omega_f-\omega_0)
=-mR(v_f-v_0).
\]

约去 \(mR\)：

\[
\frac23R(\omega_f-\omega_0)
=-(v_f-v_0).
\]

末态无滑动给出 \(\omega_f=v_f/R\)，于是

\[
\frac23(v_f-R\omega_0)=-v_f+v_0.
\]

整理：

\[
\frac53v_f=v_0+\frac23R\omega_0.
\]

所以

\[
\boxed{v_f=\frac35v_0+\frac25R\omega_0},
\]

\[
\boxed{\omega_f=\frac{v_f}{R}
=\frac35\frac{v_0}{R}+\frac25\omega_0}.
\]

冲量为

\[
\begin{aligned}
J
&=m(v_f-v_0)\\
&=\frac{2m}{5}(R\omega_0-v_0).
\end{aligned}
\]

因此初始接触点的滑动趋势 \(v_0-R\omega_0\) 决定冲量方向。

\Needspace{3\baselineskip}
\textbf{2. 初始不转 \(\omega_0=0\)}

直接代入：

\[\boxed{v_f=\frac35v_0>0},\]

\[\boxed{\omega_f=\frac{3v_0}{5R}>0}.\]

桌面给球向左冲量，使水平速度减小，同时产生上旋。

\Needspace{3\baselineskip}
\textbf{3. 初始上旋 \(\omega_0>0\)}

末态两量均为正：

\[v_f=\frac35v_0+\frac25R\omega_0>0,\]

\[\omega_f=\frac35\frac{v_0}{R}+\frac25\omega_0>0.\]

但冲量方向分三种：

\[
\begin{cases}
v_0>R\omega_0,&J<0,\quad\text{桌面对球的冲量向左};\\
v_0=R\omega_0,&J=0,\quad\text{初始接触点已无相对滑动};\\
v_0<R\omega_0,&J>0,\quad\text{桌面对球的冲量向右}.
\end{cases}
\]

无论冲量朝哪边，给定的线性组合保证 \(v_f\) 与 \(\omega_f\) 均为正。

\Needspace{3\baselineskip}
\textbf{4. 初始下旋 \(\omega_0<0\)}

写成 \(\omega_0=-|\omega_0|\)：

\[v_f=\frac35v_0-\frac25R|\omega_0|,\]

\[\omega_f=\frac35\frac{v_0}{R}-\frac25|\omega_0|.\]

由于 \(v_f=R\omega_f\)，二者符号相同。临界条件为

\[3v_0-2R|\omega_0|=0.\]

故

\[
\begin{cases}
v_0>\dfrac23R|\omega_0|,&v_f>0,\ \omega_f>0;\\[4pt]
v_0=\dfrac23R|\omega_0|,&v_f=0,\ \omega_f=0;\\[4pt]
v_0<\dfrac23R|\omega_0|,&v_f<0,\ \omega_f<0.
\end{cases}
\]

最后一种情况下球在水平方向反弹回去，并保持下旋。

\Needspace{3\baselineskip}
\textbf{能量为何不守恒}

竖直方向完全弹性不代表切向机械能守恒。题设通过``离地瞬间无滑动''指定了切向碰撞的末态约束；接触区的变形、摩擦和内部振动可以耗散平动与转动能。因此本题不能再附加切向动能守恒，否则方程过定。

\Needspace{3\baselineskip}
\textbf{一般转动惯量的公式}

若 \(I=\kappa mR^2\)，同样推导可得

\[v_f=\frac{1}{1+\kappa}v_0+\frac{\kappa}{1+\kappa}R\omega_0.\]

本题 \(\kappa=2/3\)，即得到系数 \(3/5\) 和 \(2/5\)。

\Needspace{3\baselineskip}
\textbf{一般惯量与摩擦冲量是否够大}

\(I=\kappa mR^2\) 时，统一结果为

\[v_f=\frac{v_0+\kappa R\omega_0}{1+\kappa},\quad
J=-\frac{m\kappa}{1+\kappa}(v_0-R\omega_0),\]

\[K_f-K_i=-\frac{m\kappa}{2(1+\kappa)}(v_0-R\omega_0)^2.\]

\(\kappa=2/3\) 得到正文 \(3/5,2/5\)；实心球 \(\kappa=2/5\) 则为
\(5/7,2/7\)，反弹向后的阈值是 \(R|\omega_0|>5v_0/2\)。

若碰前向下速度大小为 \(w\)，竖直完全弹性给出法向冲量
\(J_n=2mw\)。采用库仑冲量上界 \(|J|\le\mu J_n\) 时，末态无滑动至少需

\[\mu\ge\frac{\kappa|v_0-R\omega_0|}{2(1+\kappa)w}.\]

原题直接规定离桌无滑动，因而默认接触作用能提供这份冲量；不能在任意微小摩擦下照搬。解图中的
\(J\)
向右只是正方向，实际方向由初始滑动决定。题图反弹轨迹应向右上，原稿向左下的箭头已修正。

\Needspace{3\baselineskip}
\textbf{无自转与强下旋的两个回代例}

对题给空心球 \(\kappa=2/3\)，初始不转时
\(v_f=3v_0/5\)、\(R\omega_f=3v_0/5\)、\(J=-2mv_0/5\)。初始水平动能是
\(mv_0^2/2\)，末态平动与转动之和为

\[K_{t,f}=\frac12m\left(\frac{3v_0}{5}\right)^2
+\frac12\frac23mR^2\left(\frac{3v_0}{5R}\right)^2
=\frac3{10}mv_0^2.\]

少了 \(mv_0^2/5\)，与一般耗散式吻合。若 \(R\omega_0=-2v_0\)，

\[v_f=-\frac15v_0,\qquad R\omega_f=-\frac15v_0,\qquad J=-\frac65mv_0.\]

两者都反号，接触点仍无滑动。临界下旋是
\(R\omega_0=-3v_0/2\)，此时水平平动与自转都为零，而竖直分量依题设仍反弹。图示向右反弹属于无自转或上旋的情形，不能用同一个方向箭头覆盖所有分支。

\Needspace{3\baselineskip}
\textbf{从冲量消元得到一般惯量结果}

取向右为水平正向、顺时针为自转正向，\(I=\kappa mR^2\)。切向冲量 \(J\)
满足

\[m(v-v_0)=J,\qquad I(\omega-\omega_0)=-RJ.\]

若离桌瞬间达到无滑动，\(v=R\omega\)。把第二式乘 \(R/I\) 并消去
\(\omega\)，得

\[J=-\frac{m\kappa}{1+\kappa}(v_0-R\omega_0),\]
\[v=\frac{v_0+\kappa R\omega_0}{1+\kappa},\qquad
\omega=\frac{v_0+\kappa R\omega_0}{(1+\kappa)R}.\]

空心球 \(\kappa=2/3\)，便得到 \(v=3v_0/5+2R\omega_0/5\)。强下旋使
\(\omega_0<0\)，只有 \(R|\omega_0|>v_0/\kappa\) 才反向；实心球阈值为
\(5v_0/2\)。不能把``上旋''一概画成水平反弹向左，图中的反弹箭头应说明使用哪一初始旋转条件。

\Needspace{3\baselineskip}
\textbf{可实现无滑动终态还需足够摩擦冲量}

竖直入射速率为 \(w\)，题中若竖直完全弹性，法向冲量为
\(J_n=2mw\)。库仑摩擦至少要求

\[|J|\le\mu J_n\quad\Longrightarrow\quad
\mu\ge\frac{\kappa|v_0-R\omega_0|}{2(1+\kappa)w}.\]

这只是实现所需切向冲量的必要尺度条件；具体接触过程还要与题设的恢复和滑动机制一致。若不足，应在整个碰撞中保留滑动，用受限的切向冲量算出
\(v,\omega\)，而不能继续套 \(v=R\omega\)。

切向平动与自转能量的合变化为

\[\Delta K_t=-\frac{m\kappa}{2(1+\kappa)}(v_0-R\omega_0)^2\le0.\]

在 \(v_0=R\omega_0\)
时所需切向冲量为零，能量不变；无自转入射则有明确耗散。竖直速度完全恢复也不代表总机械能完全恢复，切向摩擦仍可耗能。

\newpage

\section{第 4 题}\label{ux7b2c-4-ux9898}

一块均匀薄木板原本为直板，现弯成半圆形并夹在两个静止物块之间，与两物块均为点接触。木板中性层长度为
\(\pi R\)，厚度为 \(d\)，垂直纸面方向宽度为 \(a\)，杨氏模量为
\(Y\)；木板不接触地面。求每个物块受到木板的水平作用力 \(F\)。

\begin{figure}[H]
\centering
\includegraphics[width=0.794\linewidth,height=\textheight,keepaspectratio,alt={图 4-1　弯成半圆形的薄木板及两端点接触。}]{../figures/question-04.pdf}
\caption*{图 4-1　弯成半圆形的薄木板及两端点接触。}
\end{figure}

\subsection{一、题型定位}\label{ux4e00ux9898ux578bux5b9aux4f4d-3}

先求薄板被强制弯成圆弧时的弯曲能，再用端点间距计算圆弧族的虚功。最后必须检查这种形状是否能由题给的点接触维持。

\subsection{二、知识点}\label{ux4e8cux77e5ux8bc6ux70b9-3}

\begin{itemize}
\tightlist
\item
  纯弯曲中应变沿厚度线性变化。
\item
  弹性能密度 \(u=\frac12Y\varepsilon^2\)。
\item
  截面二次矩 \(I=ad^3/12\)。
\item
  中性层长度不变条件。
\item
  广义力与势能导数。
\end{itemize}

\subsection{三、物理图像}\label{ux4e09ux7269ux7406ux56feux50cf-3}

木板弯曲后，外层被拉长、内层被压缩，中性层长度保持
\(\pi R\)。若两个物块稍向外移动，木板曲率减小、弹性能下降；木板正是通过释放这部分能量推开物块。

\subsection{四、解题入口}\label{ux56dbux89e3ux9898ux5165ux53e3-3}

设形变过程中的瞬时曲率半径为 \(\rho\)，但中性层总长度始终为

\[L=\pi R.\]

先求 \(U(\rho)\)，再求端点间距 \(s(\rho)\) 在 \(\rho=R\)
附近的变化，使用

\[F=-\frac{dU}{ds}.\]

\subsection{五、解析}\label{ux4e94ux89e3ux6790-3}

\begin{figure}[H]
\centering
\includegraphics[width=0.9\linewidth,height=\textheight,keepaspectratio,alt={图 4-2　弯曲木板厚度方向的线性应变分布。}]{../figures/solution-04.pdf}
\caption*{图 4-2　弯曲木板厚度方向的线性应变分布。}
\end{figure}

\Needspace{3\baselineskip}
\textbf{1. 厚度方向的应变}

取距中性层为 \(x\) 的纤维。中性层一小段张角为 \(d\theta\)，长度为
\(\rho d\theta\)；该纤维长度为 \((\rho+x)d\theta\)。因此相对伸长为

\[
\varepsilon(x)
=\frac{(\rho+x)d\theta-\rho d\theta}{\rho d\theta}
=\frac{x}{\rho}.
\]

弹性能密度为

\[u(x)=\frac12Y\varepsilon^2
=\frac12Y\frac{x^2}{\rho^2}.\]

\Needspace{3\baselineskip}
\textbf{2. 对截面和全长积分}

体积元为

\[dV=a\,dx\,ds.\]

所以总弹性能

\[
U=\int_0^L\int_{-d/2}^{d/2}
\frac12Y\frac{x^2}{\rho^2}a\,dx\,ds.
\]

先对厚度积分：

\[
\int_{-d/2}^{d/2}x^2dx
=\left[\frac{x^3}{3}\right]_{-d/2}^{d/2}
=\frac{d^3}{12}.
\]

于是

\[
\begin{aligned}
U(\rho)
&=\frac12Y\frac{a}{\rho^2}\frac{d^3}{12}L\\
&=\frac{Yad^3L}{24\rho^2}\\
&=\frac{\pi Yad^3R}{24\rho^2}.
\end{aligned}
\]

这里不能把 \(L\) 写成
\(\pi\rho\)，因为木板没有沿中性层伸长，恒定的是原长度 \(L=\pi R\)。

\Needspace{3\baselineskip}
\textbf{3. 端点间距对曲率半径的导数}

长度为 \(L\)、曲率半径为 \(\rho\) 的圆弧，其圆心角为

\[\Theta=\frac{L}{\rho}.\]

两端弦长为

\[s(\rho)=2\rho\sin\frac{\Theta}{2}
=2\rho\sin\frac{L}{2\rho}.\]

求导：

\[
\frac{ds}{d\rho}
=2\sin\frac{L}{2\rho}
-\frac{L}{\rho}\cos\frac{L}{2\rho}.
\]

在半圆状态 \(\rho=R\)、\(L=\pi R\) 时，

\[\frac{L}{2\rho}=\frac\pi2,\]

故

\[
\left.\frac{ds}{d\rho}\right|_{\rho=R}
=2.
\]

\Needspace{3\baselineskip}
\textbf{4. 用虚功求作用力}

弹性能导数为

\[
\frac{dU}{d\rho}
=-\frac{\pi Yad^3R}{12\rho^3}.
\]

在 \(\rho=R\)：

\[
\left.\frac{dU}{d\rho}\right|_R
=-\frac{\pi Yad^3}{12R^2}.
\]

广义力

\[
F=-\frac{dU}{ds}
=-\frac{dU/d\rho}{ds/d\rho}.
\]

代入上述两个导数：

\[
\boxed{F=\frac{\pi Yad^3}{24R^2}}.
\]

这是在额外圆弧形状约束下得到的端点间距共轭力，方向均向外；原题点接触能否维持半圆须见下面的边界条件核对。

\Needspace{3\baselineskip}
\textbf{量纲检查}

\(Y\) 的量纲为力每面积，\(ad^3/R^2\)
的量纲为面积，因此结果量纲为力。厚度以三次方出现，说明薄板抗弯刚度对厚度极敏感。

\Needspace{3\baselineskip}
\textbf{与梁弯曲公式的对应}

矩形截面二次矩 \(I=ad^3/12\)，纯弯曲能也可写为

\[U=\frac12YI\kappa^2L,\qquad\kappa=\frac1\rho.\]

代入后立即得到同一 \(U(\rho)\)。从应变密度推导一次，有助于学生理解 \(I\)
的来源。

\Needspace{3\baselineskip}
\textbf{圆弧族计算与实际点接触的区别}

上面的 \(F=\pi Yad^3/(24R^2)\)
与扫描参考解一致，但有额外条件：形变被限制在恒曲率圆弧族，并由适当分布作用维持。原题只给两端点接触，不能据此证明这个条件。

对原本直的梁，恒半径半圆要求每个截面弯矩
\(M=YI/R\)。仅有两端水平接触力时，从端点截取一段梁，其截面弯矩却是
\(M(s)=F y(s)\)，其中 \(y\) 为截面高于端点作用线的距离。在端点
\(y=0\)，弯矩应为零；半圆的恒曲率要求它仍为非零。这两项边界条件互相冲突。

因此``恰好半圆且只受端点水平力''不是普通直梁的静力平衡。严格按题面，无法给出上述数值为真实端点反力。若补充保持圆弧的形状约束，公式才是沿该约束族的等效广义力；若只保留端点夹持，则须求实际非圆形弹性曲线及其边界条件。图
4-2
的力和能量均标为圆弧约束模型，避免把形式虚功冒充已经满足局部平衡的结果。

\Needspace{3\baselineskip}
\textbf{厚度、宽度与应变的检查}

恒曲率族中最外纤维的应变绝对值为
\(d/(2\rho)\)。线弹性与薄梁计算至少要求它远小于1。因此公式对厚度的三次依赖只在
\(d\ll R\) 的范围内使用。

宽度翻倍、厚度不变时，\(I\)
与能量、形式广义力都翻倍；厚度翻倍时，它们增为8倍。若把宽度和厚度互换，弯曲方向也随之改变，不能保持同一
\(I\)。

从原长 \(L\) 的一般圆弧族出发，令 \(z=L/(2\rho)\)，端点间距导数为

\[\frac{ds}{d\rho}=2(\sin z-z\cos z).\]

形式广义力为

\[F(\rho)=\frac{YIL}{2\rho^3(\sin z-z\cos z)}.\]

半圆 \(z=\pi/2\) 时还原参考解；当 \(z\to0\) 时分母含
\(z^3/3\)，对应接近直梁的退化几何。无论取哪一个
\(z\)，这些表达式仍需外部形状约束，不能补救原题只用端点接触力的局部平衡矛盾。

\Needspace{3\baselineskip}
\textbf{从截面伸长推导弯曲能}

以中性层为 \(x=0\)，矩形截面厚度方向为 \(-d/2\le x\le d/2\)，宽度为
\(a\)。曲率为 \(1/\rho\) 时应变 \(\varepsilon=x/\rho\)。一小段原长
\(ds\) 的能量为

\[dU=\frac12Y\int_A\varepsilon^2dA\,ds
=\frac{Y}{2\rho^2}a\int_{-d/2}^{d/2}x^2dx\,ds.\]

由于 \(I=ad^3/12\)，得 \(dU=YI,ds/(2\rho^2)\)。沿固定长度 \(L=\pi R\)
的圆弧积分，

\[U(\rho)=\frac{YI L}{2\rho^2},\qquad
s(\rho)=2\rho\sin\frac{L}{2\rho}.\]

在 \(\rho=R\)，\(s'=2\)、\(U'=-YI L/R^3\)。沿强制圆弧族的虚功给
\(F=-U'/s'=\pi YI/(2R^2)\)，也就是
\(\pi Yad^3/(24R^2)\)。这里算出的是圆弧形状约束下的广义力，不能跳过下面的局部平衡检查。

\Needspace{3\baselineskip}
\textbf{为何端点力不能维持恒曲率半圆}

对从一个端点截到内部某处的梁段作自由体图。只有水平端点力时，截面弯矩等于
\(F\) 乘截面与端点之间的竖直高度差，即 \(M(s)=Fy(s)\)。端点附近
\(y\to0\)，所以 \(M\to0\)。然而从直梁弯成曲率处处为 \(1/R\)
的半圆，线弹性梁要求

\[M(s)=YI/R,\]

是非零常量。两条件在端点已经矛盾，内部也不可能同时成立。这与总能量是否写对无关，是题面边界条件本身不足。

若用沿梁分布的形状约束维持圆弧，约束反力会参与局部平衡；若端部另施加力偶，应把力偶及端部方向写入边界条件；若仅有题面的块与水平接触力，实际形状需由弹性曲线方程求出，通常不是圆。题面图继续保留原题的声称，解析图明确标出圆弧约束模型，避免把改好的图当作原题条件已经成立的证明。

\newpage

\section{第 5 题}\label{ux7b2c-5-ux9898}

竖直平行板电容器的两极板固定，间距为 \(d\)、高度为 \(h\)、面积
\(A\gg d^2\)。忽略边缘效应和空气阻力，考虑重力。

一质量为 \(M\)、电荷量为 \(q\) 的小金属球用长为 \(L\)
的绝缘细线悬挂。细线悬点固定；电容器不带电时，小球位于两板中央，高度为
\(h/2\)。

\begin{enumerate}
\def\labelenumi{\arabic{enumi}.}
\tightlist
\item
  两板电势差恒为 \(V_0\) 时，求平衡角 \(\theta_0\)；将小球从略大于
  \(\theta_0\) 的角度由静止释放，证明其作小振动并求周期。
\item
  小球通过平衡位置时剪断细线，求小球不撞击极板时 \(V_0\) 的最大值。
\item
  另一过程中，\(t=0\)
  时小球位于两板中心且静止，同时剪断细线，并加交流电压
  \(V(t)=V_0\sin\omega t\)。分别在 \(g\gg h\omega^2\) 和
  \(g\ll h\omega^2\) 两个极限下求小球不碰极板时 \(\omega\) 的范围。
\end{enumerate}

\begin{figure}[H]
\centering
\includegraphics[width=0.528\linewidth,height=\textheight,keepaspectratio,alt={图 5-1　平行板间带电摆球的受力与几何位置。}]{../figures/question-05.pdf}
\caption*{图 5-1　平行板间带电摆球的受力与几何位置。}
\end{figure}

\subsection{一、题型定位}\label{ux4e00ux9898ux578bux5b9aux4f4d-4}

恒压时电力与重力合成倾斜的等效重力；剪线后分别处理初始几何和运动。交流电压部分连续积分两次，再检查球离开板区前的最大横移。

\subsection{二、知识点}\label{ux4e8cux77e5ux8bc6ux70b9-4}

\begin{itemize}
\tightlist
\item
  平行板电场 \(E=V/d\)。
\item
  恒力场合成与等效重力。
\item
  单摆小振动。
\item
  二维匀加速运动。
\item
  正弦驱动的两次积分。
\item
  小参数与大参数渐近展开。
\end{itemize}

\subsection{三、物理图像}\label{ux4e09ux7269ux7406ux56feux50cf-4}

恒定电压时，水平电力与竖直重力合成为一条倾斜的``等效重力''。摆球围绕这条方向平衡。剪线后，水平与竖直加速度的比保持为
\(\tan\theta_0\)。交流电压从零相位启动，横向加速度虽交替变号，但在球落到底板前积累的位移并不周期性返回零。

\subsection{四、解题入口}\label{ux56dbux89e3ux9898ux5165ux53e3-4}

恒压部分先定义

\[E_0=\frac{V_0}{d},
\qquad
a_E=\frac{qV_0}{Md}.\]

交流部分先求落到底端的时间

\[t_f=\sqrt{\frac hg},\]

再将横向位移 \(x(t_f)\) 与半间距 \(d/2\) 比较。

\subsection{五、解析}\label{ux4e94ux89e3ux6790-4}

\Needspace{3\baselineskip}
\textbf{1. 平衡角}

平衡时水平方向和竖直方向分别有

\[T\sin\theta_0=qE_0,\]

\[T\cos\theta_0=Mg.\]

两式相除：

\[
\boxed{\tan\theta_0=\frac{qV_0}{Mgd}}.
\]

\Needspace{3\baselineskip}
\textbf{2. 等效重力与小振动周期}

恒定的重力和电力合成为

\[
Mg'=\sqrt{(Mg)^2+(qE_0)^2}.
\]

所以

\[
g'=\sqrt{g^2+\left(\frac{qV_0}{M d}\right)^2}.
\]

又由 \(\tan\theta_0=qE_0/(Mg)\)，

\[g'=\frac g{\cos\theta_0}.\]

令相对平衡方向的小偏角为 \(\eta=\theta-\theta_0\)。沿切向的回复力在
\(|\eta|\ll1\) 时为

\[F_t=-Mg'\sin\eta\simeq-Mg'\eta.\]

切向加速度为 \(L\ddot\eta\)，故

\[ML\ddot\eta=-Mg'\eta.\]

即

\[
\ddot\eta+\frac{g'}{L}\eta=0.
\]

因此为简谐运动，角频率和周期分别为

\[\Omega=\sqrt{\frac{g'}L},\]

\[
\boxed{
T=2\pi\sqrt{\frac{L}{g'}}
=2\pi\sqrt{\frac Lg}
\left[1+\left(\frac{qV_0}{Mgd}\right)^2\right]^{-1/4}}.
\]

\Needspace{3\baselineskip}
\textbf{3. 在平衡位置剪线后的安全条件}

``轻微抬高''表示小振幅极限；通过平衡位置时的附加速度可忽略到本问保留阶次。剪线瞬间，小球相对两板中央向电力方向偏移

\[x_0=L\sin\theta_0.\]

其离底边高度为

\[y_0=\frac h2+L(1-\cos\theta_0).\]

剪线后

\[a_x=\frac{qV_0}{Md}=g\tan\theta_0,\]

\[a_y=g.\]

从静止开始下降 \(y_0\) 时，水平附加位移与竖直位移之比等于加速度之比：

\[\Delta x=y_0\frac{a_x}{g}=y_0\tan\theta_0.\]

到达下边界高度时的总横向偏移为

\[
\begin{aligned}
x_{\rm tot}
&=L\sin\theta_0+
\left[\frac h2+L(1-\cos\theta_0)\right]\tan\theta_0\\
&=\left(L+\frac h2\right)\tan\theta_0.
\end{aligned}
\]

临界情况下恰到达极板：

\[
\left(L+\frac h2\right)\tan\theta_0=\frac d2.
\]

代入 \(\tan\theta_0=qV_0/(Mgd)\)：

\[
\left(L+\frac h2\right)\frac{qV_0}{Mgd}=\frac d2.
\]

解得

\[
\boxed{V_{0,\max}=\frac{Mgd^2}{q(2L+h)}}.
\]

\Needspace{3\baselineskip}
\textbf{4. 交流电压作用下的横向运动}

此时

\[a_x(t)=\frac{qV_0}{Md}\sin\omega t.\]

由 \(v_x(0)=0\) 积分：

\[
\begin{aligned}
v_x(t)
&=\int_0^t a_x(t')dt'\\
&=\frac{qV_0}{Md\omega}(1-\cos\omega t).
\end{aligned}
\]

再由 \(x(0)=0\) 积分：

\[
\begin{aligned}
x(t)
&=\int_0^t v_x(t')dt'\\
&=\frac{qV_0}{Md\omega}
\left(t-\frac{\sin\omega t}{\omega}\right).
\end{aligned}
\]

图 5-2
显示，即使加速度交变，\(1-\cos\omega t\ge0\)，所以速度不反向，横向位移总体单调增加。

\begin{figure}[H]
\centering
\includegraphics[width=0.823\linewidth,height=\textheight,keepaspectratio,alt={图 5-2　正弦电压从零相位启动时横向位移随时间的累积。}]{../figures/solution-05.pdf}
\caption*{图 5-2　正弦电压从零相位启动时横向位移随时间的累积。}
\end{figure}

竖直方向为自由落体。从中心高度 \(h/2\) 落到底端需要

\[
\frac h2=\frac12gt_f^2,
\]

所以

\[t_f=\sqrt{\frac hg}.\]

安全条件为

\[x(t_f)<\frac d2.\]

\Needspace{3\baselineskip}
\textbf{5. 低频极限 \(g\gg h\omega^2\)}

该条件等价于

\[\omega t_f=\omega\sqrt{\frac hg}\ll1.\]

令 \(z=\omega t_f\)，用

\[\sin z=z-\frac{z^3}{6}+O(z^5).\]

于是

\[
t_f-\frac{\sin\omega t_f}{\omega}
\simeq\frac{\omega^2t_f^3}{6}.
\]

所以

\[
x(t_f)\simeq
\frac{qV_0\omega t_f^3}{6Md}.
\]

要求 \(x(t_f)<d/2\)：

\[
\boxed{
\omega<\frac{3Md^2}{qV_0}
\left(\frac gh\right)^{3/2}}
\qquad(g\gg h\omega^2).
\]

\Needspace{3\baselineskip}
\textbf{6. 高频极限 \(g\ll h\omega^2\)}

此时 \(\omega t_f\gg1\)。括号中主项为 \(t_f\)，有

\[
x(t_f)\simeq\frac{qV_0t_f}{Md\omega}.
\]

安全条件给出

\[
\boxed{
\omega>\frac{2qV_0}{Md^2}\sqrt{\frac hg}}
\qquad(g\ll h\omega^2).
\]

\Needspace{3\baselineskip}
\textbf{渐近结果的自洽检查}

低频结果必须与 \(\omega\sqrt{h/g}\ll1\) 同时满足，高频结果必须与
\(\omega\sqrt{h/g}\gg1\)
同时满足。若代入题给参数后两个条件不兼容，说明不能使用相应渐近式，需直接求解

\[
\frac{qV_0}{Md\omega}
\left(t_f-\frac{\sin\omega t_f}{\omega}\right)=\frac d2.
\]

\Needspace{3\baselineskip}
\textbf{为何高频下位移反而变小}

电压振荡越快，每个半周期给出的速度增量越小。虽然速度始终不为负，但平均漂移速度按
\(1/\omega\) 降低，所以足够高频时小球来不及撞板。

\Needspace{3\baselineskip}
\textbf{剪线初速度的保留与零振幅极限}

从较大角度释放后第一次通过平衡，速度并非严格为零。若小振幅为
\(\eta_m\)，过平衡速率
\(u\simeq\sqrt{Lg'}\,\eta_m\)，方向为向左、向下。取向下为正，

\[x(t)=L\sin\theta_0-u\cos\theta_0\,t+\tfrac12g\tan\theta_0\,t^2,\]

\[y(t)=u\sin\theta_0\,t+\tfrac12gt^2.\]

离开下边缘的时间由 \(y(t)=y_0\) 确定，代回得到

\[x(t_f)=(L+h/2)\tan\theta_0-\frac{u t_f}{\cos\theta_0}.\]

所以正文电压阈值是 \(\eta_m\to0\)
的极限，不是任意``略高''幅度下的精确最大值；有限 \(u\)
还要检查轨迹左端最小横坐标。平衡位置须仍在板区，\(h/2+L(1-\cos\theta_0)<h\)。

若严格要求``不接触极板''，临界等号只给电压上确界，应写
\(|V_0|<Mgd^2/[|q|(2L+h)]\)；等号是到下边缘时恰触板角的边界。

交流驱动的精确安全条件为 \(\Lambda(z-\sin z)/z^2<1\)。解图采用真实的
\(x\propto z-\sin z\)，标出速度为零但不反向的
\(z=2\pi\)。余弦初相位变式则
\(x=a_E(1-\cos\omega t)/\omega^2\)，必须检查整个时间区间的峰值，不能只看离场瞬间。

\Needspace{3\baselineskip}
\textbf{零振幅阈值和严格不接触的区别}

令剪线瞬间 \(u=0\)，则 \(x\) 从初值单调增加，到底部有
\(x_f=(L+h/2)\tan\theta_0\)。不碰右板要求

\[x_f<d/2\quad\Longrightarrow\quad
\tan\theta_0<\frac d{2L+h}.\]

结合平衡条件，

\[V_0<\frac{Mgd^2}{q(2L+h)}.\]

等号对应球在下沿恰好碰右板，是不接触条件的上确界。原题说从``略大''角度释放，\(u\)
很小却不为零；因此上述值是振幅趋零的阈值，而非任意给定小振幅的精确最大电压。

若已知摆幅相对平衡为 \(\eta_m\)，小振动能量给
\(u^2\simeq Lg_{\rm eff}\eta_m^2\)。此时将 \(u\)
代回前述完整抛体轨迹，再检查两板边界。仅把末端右移减小看成``更安全''也不够，因为初始向左的速度还可能造成另一侧碰板。

\Needspace{3\baselineskip}
\textbf{剪线后的横移取决于过平衡位置的速度}

取向下为 \(y\) 正向。平衡角 \(\theta_0\) 满足
\(\tan\theta_0=qV_0/(Mgd)\)。若球从右侧经过平衡位置时速率为
\(u\)，初速度分量为 \((-u\cos\theta_0,u\sin\theta_0)\)。剪线后

\[x(t)=L\sin\theta_0-u\cos\theta_0t+\frac12g\tan\theta_0t^2,\]
\[y(t)=u\sin\theta_0t+\frac12gt^2.\]

到底部的正根满足 \(y(t_f)=h/2\)，而

\[x(t_f)=(L+h/2)\tan\theta_0-\frac{ut_f}{\cos\theta_0}.\]

最大横坐标发生在凸函数区间的端点；最小横坐标可能发生在
\(t_m=u\cos\theta_0/(g\tan\theta_0)\)，要看它是否落在
\([0,t_f]\)。因此有限振幅下应检查 \(x(0),x(t_f)\) 及可达的
\(x(t_m)\)，并同时保证两侧边界
\(-d/2<x<d/2\)。只检查落地横坐标会漏掉向左先移的一段。

\Needspace{3\baselineskip}
\textbf{交流部分的两个频率极限}

定义 \(a=qV_0/(Md)\)，初位移、初速度为零。由 \(\ddot x=a\sin\omega t\)
连续积分得到

\[\dot x=\frac a\omega(1-\cos\omega t)\ge0,
\qquad x=\frac a{\omega^2}(\omega t-\sin\omega t).\]

从高度 \(h/2\) 落到底部，\(t_f=\sqrt{h/g}\)。由于 \(x\) 不减，只需检查
\(x(t_f)<d/2\)。令 \(z=\omega t_f\)，精确不等式为

\[\frac a{\omega^2}(z-\sin z)<\frac d2.\]

低频 \(z\ll1\) 时，\(z-\sin z\sim z^3/6\)，于是
\(\omega<3Md^2(g/h)^{3/2}/(qV_0)\)。高频 \(z\gg1\)
时，有界的正弦只给相对 \(O(1/z)\) 修正，得到
\(\omega>2qV_0\sqrt{h/g}/(Md^2)\)。得到的频率还须落在各自假设范围，不能把两条渐近式当成中间频率的精确分界。

\Needspace{3\baselineskip}
\textbf{交流变式：有初始横向速度时须检查内点}

令 \(A=qV_0/(Md\omega)>0\)，起点在两板中线，初始横速为
\(u_0\)。正弦电压下

\[x(t)=u_0t+A\left(t-\frac{\sin\omega t}{\omega}\right),\]
\[v_x(t)=u_0+A(1-\cos\omega t).\]

若 \(u_0\ge0\)，仍单调向右；若
\(u_0\le-2A\)，则始终向左。这两种情形只需比较终点横移与半板间距。若
\(-2A<u_0<0\)，速度可能变号，所有内部极值时刻满足

\[\cos\omega t=1+\frac{u_0}{A}.\]

把落在 \(0<t<t_f\) 的这些时刻与0、\(t_f\) 一起代入 \(x(t)\)，要求各处
\(|x|<d/2\)。这一有限候选点检查同时覆盖两块板；仅检查终点可能漏掉``先碰左板后返回''的轨迹。

若换成余弦电压并从静止开始，则

\[x(t)=\frac{qV_0}{Md\omega^2}(1-\cos\omega t).\]

它没有正弦初相下的线性漂移，但也并非一直不减。若
\(t_f\ge\pi/\omega\)，最大横移为

\[x_{\max}=\frac{2qV_0}{Md\omega^2}.\]

第一峰以后会返回。图 5-2
的单调曲线仅对应原题正弦初相，不能沿用作这个变式的图。

\newpage

\section{第 6 题}\label{ux7b2c-6-ux9898}

用性质均匀、单位长度电阻恒定的电阻丝构造一个无限递归三角网络。最外层是边长为
\(l\) 的正三角形；在每个角落继续生成边长为上一层 \(1/3\)
的相似三角形，三个角落完全相同。长度为 \(l\) 的电阻丝电阻为 \(R_0\)。

\begin{enumerate}
\def\labelenumi{\arabic{enumi}.}
\tightlist
\item
  求顶点 \(A\) 与左下顶点 \(B\) 间的等效电阻；
\item
  空间中有均匀变化的垂直磁场
  \(\partial B/\partial t=k_0\)。定义最外层三角形磁通变化率
\end{enumerate}

\[\mathcal E_0=k_0\frac{\sqrt3}{4}l^2.\]

取沿最外边 \(AB\) 的路径，求从 \(A\) 到 \(B\) 的电阻压降线积分
\(U_{AB}=\sum I_ir_i\)； 3. 求整个网络的焦耳热功率。

\begin{figure}[H]
\centering
\includegraphics[width=0.707\linewidth,height=\textheight,keepaspectratio,alt={图 6-1　三角形角落按比例 1/3 无限递归的电阻网络。}]{../figures/question-06.pdf}
\caption*{图 6-1　三角形角落按比例 \(1/3\) 无限递归的电阻网络。}
\end{figure}

\subsection{一、题型定位}\label{ux4e00ux9898ux578bux5b9aux4f4d-5}

把三个角落当作缩小的三端网络，可求电阻不动点。感应问题还须固定路径与磁通归一化；用有限层节点计算核对环流和总功率。

\subsection{二、知识点}\label{ux4e8cux77e5ux8bc6ux70b9-5}

\begin{itemize}
\tightlist
\item
  电阻与长度成正比。
\item
  三角形---星形变换。
\item
  三端口对称网络的等效参数。
\item
  法拉第电磁感应定律。
\item
  线性网络叠加原理。
\item
  焦耳功率的二次型。
\end{itemize}

\subsection{三、物理图像}\label{ux4e09ux7269ux7406ux56feux50cf-5}

把网络放大或缩小后，角落中的电路仍与整体同形。稳恒电源问题中可用端点电压；但时变磁场中电场非保守，不同路径的电阻压降线积分可以不同，不能把
\(U_{AB}\) 简单称为端点电势差。

\subsection{四、解题入口}\label{ux56dbux89e3ux9898ux5165ux53e3-5}

把无限网络看成一个关于 \(A,B,C\) 完全对称的三端黑箱。先用三角形边参数
\(xR_0\) 表示这个黑箱，再把它变换成三臂均为 \(xR_0/3\)
的星形网络。这样，本级连接段和缩小副本的贡献都能在同一种参数下比较，最后再换算成可测量的两端电阻
\(R_{AB}\)。

\subsection{五、解析}\label{ux4e94ux89e3ux6790-5}

\begin{figure}[H]
\centering
\includegraphics[width=0.587\linewidth,height=\textheight,keepaspectratio,alt={图 6-2　利用三端口对称性表示每个角落的自相似子网络。}]{../figures/solution-06.pdf}
\caption*{图 6-2　利用三端口对称性表示每个角落的自相似子网络。}
\end{figure}

\Needspace{3\baselineskip}
\textbf{1. 等效电阻的自相似参数}

设完整无限网络的三角形等效电路中，每一条边的电阻为

\[r_\triangle=xR_0.\]

这里的 \(x\) 是``三角形边参数''，还不是 \(A\)、\(B\)
两端看到的等效电阻系数。把等边三角形变换成等臂星形。三角形---星形变换公式为

\[
r_A=\frac{r_{AB}r_{AC}}{r_{AB}+r_{BC}+r_{CA}}.
\]

三条边均为 \(xR_0\)，故每条星臂都是

\[
r_\star=\frac{(xR_0)^2}{3xR_0}=\frac{xR_0}{3}.
\]

现在看递归图中的一个角落。子图线性尺寸缩小为原来的
\(1/3\)，所以其中每段电阻也缩小为原来的
\(1/3\)。把三个角落子网络依次作三角形---星形变换，并利用三条中线的镜面对称合并等势节点后，大网络的一条等效星臂由三部分组成：

\[
\underbrace{\frac13\left(\frac{R_0}{3}+\frac{2xR_0}{9}\right)}_{\text{本级连接段与两侧子网络经对称分流后的贡献}}
+
\underbrace{\frac{xR_0}{9}}_{\text{正对该顶点的子网络星臂}}.
\]

而完整网络本身的一条星臂又必须等于 \(xR_0/3\)。因此自相似不动点方程是

\[
\frac13\left(\frac{R_0}{3}+\frac{2xR_0}{9}\right)
+\frac{xR_0}{9}
=\frac{xR_0}{3}.
\]

先展开左边：

\[
\frac{R_0}{9}+\frac{2xR_0}{27}+\frac{3xR_0}{27}
=\frac{xR_0}{3}.
\]

合并同类项，并把右边通分为分母 \(27\)：

\[
\frac{3R_0}{27}+\frac{5xR_0}{27}
=\frac{9xR_0}{27}.
\]

两边同乘 \(27/R_0\)：

\[3+5x=9x.\]

所以

\[4x=3,\]

\[\boxed{x=\frac34}.\]

\(A\)、\(B\) 两端之间，一条路径直接通过三角形的一条边
\(xR_0\)，另一条路径通过另外两条边串联，电阻为
\(2xR_0\)。两条路径并联，故

\[
\begin{aligned}
R_{AB}
&=xR_0\parallel 2xR_0\\
&=\frac{xR_0\cdot2xR_0}{xR_0+2xR_0}\\
&=\frac23xR_0\\
&=\frac23\times\frac34R_0\\
&=\frac12R_0.
\end{aligned}
\]

因此

\[\boxed{R_{AB}=\frac12R_0}.\]

\Needspace{3\baselineskip}
\textbf{2. 感应电动势的面积缩放}

最外三角形面积为

\[S_0=\frac{\sqrt3}{4}l^2,\]

其磁通变化率大小为

\[\mathcal E_0=k_0S_0.\]

任一缩小比为 \(1/3\) 的子三角形面积为原来的
\(1/9\)，故其环路感应电动势为

\[\mathcal E_1=\frac{\mathcal E_0}{9}.\]

网络关于三条中线完全对称。先把端口电流也保留下来。若按图示方向定义两个独立端口电流
\(I_1,I_2\)，线性叠加给出

\[
U_{AB}=\frac{\mathcal E_0}{3}
+\frac{R_0}{2}I_1+\frac{R_0}{4}I_2.
\]

式中三个系数可分别由三个独立工况确定：只保留磁场变化时得到 \(1/3\)；令
\(\mathcal E_0=I_2=0\) 时，\(U_{AB}=I_1R_{AB}=I_1R_0/2\)；令
\(\mathcal E_0=I_1=0\) 时，由三端口对称分流得到 \(U_{AB}=I_2R_0/4\)。

本问外端口开路，故 \(I_1=I_2=0\)。于是

\[
\boxed{U_{AB}=\frac13\mathcal E_0}.
\]

必须强调，\(U_{AB}=\sum I_ir_i\) 是指定路径上的电阻压降和。在时变磁场中

\[\oint\bm E\cdot d\bm l\ne0,\]

所以不能把它解释为与路径无关的静电势差。

\Needspace{3\baselineskip}
\textbf{3. 功率二次型及系数递推}

网络由线性电阻组成，所有支路电流都对外加端口电流和环路电动势呈线性关系。因此总焦耳功率必是这些源参数的二次齐次式。取两个独立端口电流为
\(I_1,I_2\)，最外层三角形的感应电动势大小为 \(\mathcal E\)，可写成

\[
P=B_1I_1^2+B_2I_1I_2+B_3I_2^2
+\frac{\beta}{R_0}\mathcal E^2+B_5\mathcal EI_1+B_6\mathcal EI_2.
\]

这里把纯感应项写成 \(\beta\mathcal E^2/R_0\)，使 \(\beta\)
成为无量纲常数。这个写法在缩小子网络中尤其重要：子网络的基准电阻变成
\(R_0/3\)，不能仍把分母写成 \(R_0\)。

先令 \(\mathcal E=0\)。利用第1问
\(R_{AB}=R_0/2\)，分别考查只从一个端口通入电流、交换两个端口以及同时通入两路电流的三种工况，可得

\[B_1=B_3=\frac{R_0}{2},\qquad B_2=\frac{R_0}{2}.\]

再利用镜面对称：把磁场变化方向反转而保持端口电流不变，关于对称轴成对支路的交叉功率互相抵消，故

\[B_5=B_6=0.\]

于是功率二次型化为

\[
P=\frac{R_0}{2}I_1^2+\frac{R_0}{2}I_1I_2
+\frac{R_0}{2}I_2^2+\frac{\beta}{R_0}\mathcal E^2.
\]

下面只保留感应源，即令 \(I_1=I_2=0\)，此时整个网络的功率为

\[P=\frac{\beta\mathcal E_0^2}{R_0}.\]

把三个角落子图化为两端口。它们接在中央大环上的两端电阻为
\(R_0/6\)；本级连接段各为 \(R_0/3\)，所以中央等效环电阻是
\(3(R_0/6+R_0/3)=3R_0/2\)。

中央六边形面积为外三角的 \(2/3\)。每个角落的感应环流对这条环路再贡献
\(\mathcal E_0/27\)，三处合计 \(\mathcal E_0/9\)。等效驱动为
\((2/3+1/9)\mathcal E_0=7\mathcal E_0/9\)，故连接段环流

\[I_t=\frac{7\mathcal E_0/9}{3R_0/2}
=\boxed{\frac{14\mathcal E_0}{27R_0}}.\]

图 6-2 同时标出连接段和子网络电阻，两个来源的电动势不可漏掉。采用题面
\(\mathcal E_0=k_0\sqrt3l^2/4\) 的归一化，原稿 \(7\mathcal E_0/(54R_0)\)
少了因子 \(4\)。

每个角落的端口驱动功率为 \(R_0 I_t^2/6\)，自身感应功率为
\(\beta\mathcal E_0^2/(27R_0)\)。三条连接段总功率为 \(R_0I_t^2\)，因此

\[\frac{\beta\mathcal E_0^2}{R_0}
=\frac{\beta\mathcal E_0^2}{9R_0}+\frac32R_0I_t^2.\]

代入环流后

\[\beta=\frac98\frac32\left(\frac{14}{27}\right)^2
=\frac{49}{108},\]

\[\boxed{P=\frac{49\mathcal E_0^2}{108R_0}
=\frac{49k_0^2l^4}{576R_0}}.\]

这是对原稿功率系数的修正，较其 \(49/1728\) 大 \(16\)
倍。独立有限层模型直接按真实节点及 \(\int\bm E\cdot d\bm l\)
建立支路方程，不借用不动点答案：第 0、1、2、3、4、5 层的
\(PR_0/\mathcal E_0^2\) 依次约为
\(0.333333,0.400000,0.424289,0.437064,0.444315,0.448437\)，趋向
\(49/108\simeq0.453704\)；两端电阻也从 \(2R_0/3\) 收敛到 \(R_0/2\)。

\Needspace{3\baselineskip}
\textbf{关于第1问参数的统一}

递归图中可以选择``子网络外三角边电阻''\,``三角---星形变换后的星臂电阻''或``任意两外顶点的等效电阻''为未知量。三种记号的不动点数值不同。正文选用三角形边参数
\(x=3/4\)，随后再由 \(R_{AB}=2xR_0/3\) 得到
\(R_{AB}=R_0/2\)，两者不能混写成同一个未知量。

\Needspace{3\baselineskip}
\textbf{为什么功率收敛}

第 \(n\) 级三角形数量按 \(3^n\) 增长，面积感应电动势按 \(9^{-n}\)
下降，电阻按 \(3^{-n}\) 下降。单元功率数量级为
\(\mathcal E_n^2/R_n\propto 27^{-n}\)，乘以数量 \(3^n\) 后为
\(9^{-n}\)，因此级数绝对收敛。

\Needspace{3\baselineskip}
\textbf{先由电阻递推检查端口支路}

把一个三角子网络的等效三角形边电阻记为 \(xR_0\)，这里端口两点电阻是
\(2xR_0/3\)。将每个子三角作 \(\Delta-Y\) 变换后，三分缩放和三条连接段给

\[x_{d+1}=\frac13+\frac59x_d,\qquad x_0=1.\]

解线性递推，

\[x_d=\frac34+\frac14\left(\frac59\right)^d,\]
\[R_{AB}^{(d)}=R_0\left[\frac12+\frac16\left(\frac59\right)^d\right].\]

第0层是普通三角形的 \(2R_0/3\)，第1层为 \(16R_0/27\)，无限层趋
\(R_0/2\)。子图缩为1/3后，其两端口电阻趋
\(R_0/6\)，这就是中央感应环里使用的子图电阻。

三角等效边 \(3R_0/4\)、两端电阻 \(R_0/2\)、星形每臂 \(R_0/4\)
是三个不同数；图中必须说明是哪一种，否则后续功率的连接电阻会差因子。

\Needspace{3\baselineskip}
\textbf{有限层图直接建立节点方程}

把外边长及单位长度电阻归一为1，外三角面积为
\(\sqrt3/4\)。每次将三个角落缩为原来的三分之一，并保留其间连接段；图 6-1
的每一条线都对应实际支路，中心没有凭空添加节点。

支路 \(i\to j\) 的电阻是长度
\(R_{ij}=|\bm r_j-\bm r_i|\)。对向纸内增加的磁场，可选择

\[\bm E(\bm r)=\frac{k_0}{2}(-y,x),\]

因此直线支路的动生源可由线积分直接算出：

\[e_{ij}=\frac{k_0}{2}(x_i y_j-y_i x_j).\]

选定统一方向后，写 \(I_{ij}=(V_i-V_j+e_{ij})/R_{ij}\)，各节点满足
\(\sum_j I_{ij}=0\)。只需固定一个节点电势消除规范自由度，不需要事先假定功率系数。解出电流后用
\(P=\sum R_{ij}I_{ij}^2\)，即可独立检查解析的不动点。

对于两端电阻核验，关闭感应源，在 \(A,B\)
节点注入和抽出单位电流，再由节点电压差求等效电阻。两种节点方程使用同一真实拓扑，避免用错误的中央星形图互相``验证''。

\Needspace{3\baselineskip}
\textbf{自相似缩放时三个量的幂次不同}

子图长度缩为1/3，电阻缩为1/3，面积与感应电动势缩为1/9。因此子图自身感应功率缩为

\[\frac{(1/9)^2}{1/3}=\frac1{27}.\]

三个子图合计是本级功率的1/9，而非1/3或1/27。由端口电阻 \(R_0/6\)
和连接段 \(R_0/3\)，中央环的外驱动耗散是
\(3(R_0/6+R_0/3)I_t^2=3R_0I_t^2/2\)。代入
\(I_t=14\mathcal E_0/(27R_0)\)，

\[P(1-1/9)=\frac32R_0\left(\frac{14\mathcal E_0}{27R_0}\right)^2,\]
\[P=\frac{49\mathcal E_0^2}{108R_0}.\]

功率与 \(k_0^2\)、\(l^4\) 成正比，与 \(R_0\) 成反比；这里 \(R_0\)
是外边总电阻，若改用单位长度电阻 \(\lambda\)，还应代
\(R_0=\lambda l\)，于是总体随长度变为
\(l^3\)。改变归一化后不做这一替换，会把同一结果误认为不同的尺寸律。

\Needspace{3\baselineskip}
\textbf{最低层节点模型的手算例}

没有递归时，三角形三个顶点和三条电阻均为 \(R_0\)。每条边电动势可取
\(\mathcal E_0/3\)，三个节点电势可取相等，环流为

\[I_0=\frac{\mathcal E_0}{3R_0}.\]

总功率是

\[P_0=3R_0I_0^2=\frac{\mathcal E_0^2}{3R_0}.\]

将外驱动关闭，A、B之间一条 \(R_0\) 与另一条 \(2R_0\) 并联，得
\(R_{AB}^{(0)}=2R_0/3\)。这是数值节点脚本第0层应恢复的两个结果。

第1层的三角子图边长为1/3，每个小环自感应电动势为
\(\mathcal E_0/9\)，不是 \(\mathcal E_0/3\)。三条大连接边电阻均
\(R_0/3\)。端口解的两点电阻为 \(16R_0/27\)，所有支路感应功率为
\(2\mathcal E_0^2/(5R_0)\)。它大于第0层的1/3，但小于无限层的49/108；新增子回路会改变电流分配，不能把第0层环流不变地复制到全部小环。

节点法中电势只负责使支路电流满足守恒，电动势仍由非保守场的线积分给出。某条边的电压与该边源不能混作同一变量；外边AB的测量路径约定也需要保留。

\newpage

\section{第 7 题}\label{ux7b2c-7-ux9898}

\begin{enumerate}
\def\labelenumi{\arabic{enumi}.}
\tightlist
\item
  体积为 \(V\)
  的房间由绝热材料筑成，气密性良好，仅有一条微小缝隙与外界相通，外界压强恒定。空气视为摩尔质量
  \(M_a\)、绝热指数 \(\gamma=1.4\) 的理想气体；温度为 \(T_0\)
  时房内空气密度为 \(\rho\)。求把房间缓慢加热到 \(T_1\) 所需热量。
\item
  两个与外界绝热的物体初温分别为 \(T_a\)、\(T_b\)，且
  \(T_a>T_b\)；恒定热容分别为
  \(C_a\)、\(C_b\)。用可逆卡诺热机连接两物体并对外做功，直至达到共同温度。求末温和最大做功。
\end{enumerate}

\begin{figure}[H]
\centering
\includegraphics[width=0.892\linewidth,height=\textheight,keepaspectratio,alt={图 7-1　定压泄气房间与由可逆热机连接的有限热源、冷源。}]{../figures/question-07.pdf}
\caption*{图 7-1　定压泄气房间与由可逆热机连接的有限热源、冷源。}
\end{figure}

\subsection{一、题型定位}\label{ux4e00ux9898ux578bux5b9aux4f4d-6}

房间升温时会泄气，热量要计流出焓；两有限热源经可逆热机达到共同温度时，则先用总熵不变求末温，再由能量差求功。

\subsection{二、知识点}\label{ux4e8cux77e5ux8bc6ux70b9-6}

\begin{itemize}
\tightlist
\item
  理想气体状态方程。
\item
  定压开放系统的焓流。
\item
  空气 \(C_{p,m}=\gamma R/(\gamma-1)=7R/2\)。
\item
  可逆卡诺过程与熵守恒。
\item
  有限热源最大功。
\end{itemize}

\subsection{三、物理图像}\label{ux4e09ux7269ux7406ux56feux50cf-6}

房间体积固定但与恒压外界相通。升温时房内物质的量减少，流出的空气带走焓；房内总内能实际上保持不变，输入热量用于补偿流出焓。第二问中热机不断在变化的高、低温之间可逆工作，最终共同温度由两物体总熵不变决定。

\subsection{四、解题入口}\label{ux56dbux89e3ux9898ux5165ux53e3-6}

第1问先由定压、定容条件写

\[n(T)=\frac{pV}{RT}.\]

第2问写可逆微元熵平衡

\[\frac{C_a,dT_a}{T_a}+
\frac{C_b,dT_b}{T_b}=0.\]

\subsection{五、解析}\label{ux4e94ux89e3ux6790-6}

\Needspace{3\baselineskip}
\textbf{1. 房间内物质的量随温度变化}

初态状态方程

\[pV=n_0RT_0.\]

由密度

\[n_0=\frac{\rho V}{M_a},\]

所以外界恒压为

\[
p=\frac{n_0RT_0}{V}
=\frac{\rho RT_0}{M_a}.
\]

升温过程中房内始终与外界机械平衡，故压强仍为 \(p\)。任意温度 \(T\) 时

\[n(T)=\frac{pV}{RT}.\]

\Needspace{3\baselineskip}
\textbf{2. 开放系统微元热量}

升高 \(dT\) 时，房内物质的量变化

\[
dn=-\frac{pV}{RT^2}dT<0.
\]

流出的物质的量为 \(-dn\)，每摩尔空气携带焓
\(C_{p,m}T\)。房间内空气总内能为

\[
U=nC_{V,m}T
=\frac{pV}{RT}C_{V,m}T
=\frac{pVC_{V,m}}{R},
\]

它在整个过程中保持不变。因此输入热量恰好补偿流出焓：

\[
dQ=C_{p,m}T(-dn).
\]

代入 \(dn\)：

\[
dQ=C_{p,m}\frac{pV}{RT}dT.
\]

对空气 \(\gamma=1.4=7/5\)，

\[
C_{p,m}=\frac{\gamma}{\gamma-1}R
=\frac72R.
\]

所以

\[
dQ=\frac72\frac{pV}{T}dT.
\]

从 \(T_0\) 积分到 \(T_1\)：

\[
\begin{aligned}
Q
&=\frac72pV\int_{T_0}^{T_1}\frac{dT}{T}\\
&=\frac72pV\ln\frac{T_1}{T_0}.
\end{aligned}
\]

代入 \(p=\rho RT_0/M_a\)：

\[
\boxed{
Q=\frac{7\rho VRT_0}{2M_a}
\ln\frac{T_1}{T_0}}.
\]

\Needspace{3\baselineskip}
\textbf{3. 两有限物体的可逆熵平衡}

设最终共同温度为
\(T_f\)。热机及两物体与外界绝热，且过程可逆，所以两物体总熵变化为零：

\[
\Delta S_a+\Delta S_b=0.
\]

恒定热容物体从 \(T_i\) 变到 \(T_f\) 的熵变为

\[\Delta S=C\ln\frac{T_f}{T_i}.\]

故

\[
C_a\ln\frac{T_f}{T_a}
+C_b\ln\frac{T_f}{T_b}=0.
\]

展开对数：

\[
(C_a+C_b)\ln T_f
=C_a\ln T_a+C_b\ln T_b.
\]

因此

\[
\boxed{
T_f=T_a^{\frac{C_a}{C_a+C_b}}
T_b^{\frac{C_b}{C_a+C_b}}}.
\]

图 7-2 表示末温是以热容为权重的几何平均，而不是算术平均。

\begin{figure}[H]
\centering
\includegraphics[width=0.732\linewidth,height=\textheight,keepaspectratio,alt={图 7-2　可逆连接过程中两物体总熵守恒所确定的共同温度。}]{../figures/solution-07.pdf}
\caption*{图 7-2　可逆连接过程中两物体总熵守恒所确定的共同温度。}
\end{figure}

\Needspace{3\baselineskip}
\textbf{4. 最大做功}

两物体不对外作体积功，内能变化可用热容乘温差表示。初总能量为

\[U_i=C_aT_a+C_bT_b+\text{常数},\]

末总能量为

\[U_f=(C_a+C_b)T_f+\text{常数}.\]

热机对外最大做功等于两物体总能量的减少：

\[
\boxed{
W_{\max}=C_aT_a+C_bT_b-(C_a+C_b)T_f}.
\]

代入末温可写成

\[
\boxed{
W_{\max}=C_aT_a+C_bT_b
-(C_a+C_b)
T_a^{\frac{C_a}{C_a+C_b}}
T_b^{\frac{C_b}{C_a+C_b}}}.
\]

\Needspace{3\baselineskip}
\textbf{第1问的控制体积核对}

开放系统第一定律可写为

\[\delta Q=dU+h_{\rm out}(-dn).\]

本题
\(dU=0\)，\(h_{\rm out}=C_pT\)，立即得到正文结果。若把房间错误地当成封闭系统，就会算成
\(n_0C_V(T_1-T_0)\) 或 \(n_0C_p(T_1-T_0)\)，两者都没有反映物质的量随
\(1/T\) 减少。

\Needspace{3\baselineskip}
\textbf{最大功与不可逆混合的比较}

若直接让两个物体接触，末温为热容加权算术平均，不能对外做功且总熵增加。使用可逆热机时末温为加权几何平均，低于算术平均，差出的能量转化为有用功。

\Needspace{3\baselineskip}
\textbf{两个可复算的数值与变式}

等热容 \(C\)、初温 \(400\,\mathrm K\) 和 \(300\,\mathrm K\) 时，可逆末温
\(\sqrt{120000}\simeq346.410\,\mathrm K\)，最大功
\(C(700-2\sqrt{120000})\simeq7.180C\)。直接接触末温
\(350\,\mathrm K\)，熵增
\(C\ln(350^2/120000)>0\)。解图用冷热温度的实际对数关系标出共同末点，不再用一条任意
\(S(T)\) 曲线代替两物体总熵约束。

若两物体热容为 \(C_a(T)=a_aT\)、\(C_b(T)=a_bT\)，熵变分别为
\(a_a(T_f-T_a)\) 和 \(a_b(T_f-T_b)\)，故

\[T_f=\frac{a_aT_a+a_bT_b}{a_a+a_b},\qquad
W_{\max}=\tfrac12[a_aT_a^2+a_bT_b^2-(a_a+a_b)T_f^2].\]

这里末温是以 \(a\)
为权重的算术平均；不能把恒热容的几何平均不加改变地照搬。

\Needspace{3\baselineskip}
\textbf{小温差时可逆输出功为二阶量}

等热容 \(C\)，令初温分别为
\(T+\delta,T-\delta\)，\(|\delta|\ll T\)。可逆终温与功为

\[T_f=\sqrt{T^2-\delta^2}
=T-\frac{\delta^2}{2T}+O(\delta^4/T^3),\] \[W_{\max}=2C(T-T_f)
=\frac{C\delta^2}{T}+O(C\delta^4/T^3).\]

虽然热交换量为温差的一阶量，可输出功却是二阶量，因为卡诺效率本身也随温差变小。若直接以初始效率乘全过程总热量，就把变化效率当常量，不能得到正确结果。

直接混合的终温为 \(T\)，熵增加

\[\Delta S=C\ln\frac{T^2}{T^2-\delta^2}
\simeq\frac{C\delta^2}{T^2}>0.\]

比较两过程可以解释为何可逆热机终温略低：热量的一部分已转成输出功。这一比较假设两个物体的热容为常量，不含相变和外界漏热。

\Needspace{3\baselineskip}
\textbf{开放房间不能采用封闭定压加热公式}

房间体积固定、压强固定，升温时 \(nT=pV/R\)
恒定。虽然每摩尔的内能上升，房内摩尔数下降，使总内能 \(nC_VT\)
保持不变。流出量

\[-dn=\frac{pV}{RT^2}dT\]

带走焓 \(C_pT(-dn)\)，包含随流体跨越边界的流动功。因此输入热量不是
\(n_0C_p(T_1-T_0)\)，而是

\[Q=\int_{T_0}^{T_1}C_pT(-dn)=\frac{C_p}{R}pV\ln\frac{T_1}{T_0}.\]

小温升 \(\Delta T/T_0\ll1\) 时，展开对数得到
\(Q\simeq n_0C_p\Delta T\)，所以封闭定压公式在一阶小温升上碰巧相同；二阶开始不同，不能凭这个极限把两个模型混同。若门窗关闭且空气质量固定，刚性房间则应使用
\(n_0C_V\Delta T\)，压强会随温度上升。

\Needspace{3\baselineskip}
\textbf{有限热源的最大功由熵约束决定}

对恒定热容的两个物体，总熵不变给

\[C_a\ln T_a+C_b\ln T_b=(C_a+C_b)\ln T_f.\]

所以 \(T_f\) 是按热容权重的几何平均。总内能下降就是最大可逆输出功：

\[W_{
\max}=C_aT_a+C_bT_b-(C_a+C_b)T_f.\]

加权算术平均不小于几何平均，保证
\(W_{\max}\ge0\)，等温时为零。若两物体直接接触而不做功，终温为加权算术平均；相同初态下，它高于可逆热机终温，多出的内能正是可逆过程本可输出的功。

等热容时，初始400 K和300 K给 \(T_f=\sqrt{120000}\) K。图 7-2
以冷热物体温度为两轴，可逆轨迹 \(T_aT_b=120000\) 与对角线 \(T_a=T_b\)
的交点就是终态。它不是单个物体熵随温度的曲线。

\Needspace{3\baselineskip}
\textbf{不等热容例子检查几何平均权重}

取 \(C_a=2C,C_b=C\)、初温 \(T_a=400\) K、\(T_b=300\) K。熵平衡为

\[2\ln(T_f/400)+\ln(T_f/300)=0,\] \[T_f^3=400^2\times300,\]

所以 \(T_f\simeq363.4\) K，最大功为

\[W_{
\max}=C(2\times400+300-3T_f)\simeq9.7C\ \mathrm K.\]

直接混合的终温为 \((2\times400+300)/3\simeq366.7\)
K，比可逆热机终温高。热容较大的热物体权重为2/3，终温更接近400
K；把几何平均权重写反，会得到更接近300 K的温度，违背热容账目。

若热容随温度变化，熵应写 \(\int C(T)dT/T\)，内能变化写
\(\int C(T)dT\)，两者不能用同一个常数因子代替。遇到相变还需加潜热熵变，原来的幂乘积公式不再适用。

\newpage

\section{第 8 题}\label{ux7b2c-8-ux9898}

把每个原子视为可极化介质小球，其感生电偶极矩满足
\(\bm p=\alpha\bm E\)。一个双原子分子由两个质量均为 \(m\)
的小球和一根劲度系数为 \(k\)、原长为 \(R\)
的轻弹簧组成。忽略一个原子产生的电场对另一个原子极化率的反馈，且形变量远小于
\(R\)。

\begin{enumerate}
\def\labelenumi{\arabic{enumi}.}
\tightlist
\item
  当 \(k\to\infty\)，沿原子连线施加恒定外电场
  \(E\)，求两原子间电磁作用力大小。
\item
  \(k\) 很大但有限，沿原子连线施加
  \(E(t)=E_0\cos\omega t\)。求稳态时原子间距改变量 \(\Delta(t)\)
  的平均值、振幅和角频率。
\item
  若外电场垂直于原子连线，结果如何？
\end{enumerate}

\begin{figure}[H]
\centering
\includegraphics[width=0.9\linewidth,height=\textheight,keepaspectratio,alt={图 8-1　可极化双原子模型及连接两原子的弹簧。}]{../figures/question-08.pdf}
\caption*{图 8-1　可极化双原子模型及连接两原子的弹簧。}
\end{figure}

\subsection{一、题型定位}\label{ux4e00ux9898ux578bux5b9aux4f4d-7}

两偶极沿分子轴时相吸，并排时相斥。外场正弦变化使作用力含平均项和二倍频项，用原子间距作坐标时惯性系数为约化质量。

\subsection{二、知识点}\label{ux4e8cux77e5ux8bc6ux70b9-7}

\begin{itemize}
\tightlist
\item
  感生电偶极矩 \(p=\alpha E\)。
\item
  偶极子的轴线上和赤道面电场。
\item
  偶极相互作用能与径向力。
\item
  两体相对坐标和约化质量。
\item
  含恒力项的受迫振动。
\item
  二倍频驱动和共振条件。
\end{itemize}

\subsection{三、物理图像}\label{ux4e09ux7269ux7406ux56feux50cf-7}

沿分子轴的外场使两偶极矩同轴同向，两个原子相吸，弹簧平均压缩；垂直外场使偶极矩并列同向，两个原子相斥，弹簧平均伸长。因为
\(p\propto\cos\omega t\) 而偶极力 \(F\propto p^2\)，原子间距以
\(2\omega\) 而不是 \(\omega\) 振动。

\subsection{四、解题入口}\label{ux56dbux89e3ux9898ux5165ux53e3-7}

统一用相对坐标

\[\Delta=r-R,\]

约化质量

\[\mu=\frac m2.\]

先求偶极相互作用力 \(F_{dd}(t)\)，再写

\[\mu\ddot\Delta+k\Delta=F_{dd}(t).\]

\subsection{五、解析}\label{ux4e94ux89e3ux6790-7}

\begin{figure}[H]
\centering
\includegraphics[width=0.755\linewidth,height=\textheight,keepaspectratio,alt={图 8-2　平行外场下同轴偶极相吸，垂直外场下并列偶极相斥。}]{../figures/solution-08.pdf}
\caption*{图 8-2　平行外场下同轴偶极相吸，垂直外场下并列偶极相斥。}
\end{figure}

\Needspace{3\baselineskip}
\textbf{1. 平行外场下的偶极作用}

每个原子的偶极矩大小为

\[p=\alpha E.\]

一个偶极子在轴线上距其 \(r\) 处产生的电场大小为

\[E_{\rm ax}=\frac{1}{4\pi\varepsilon_0}\frac{2p}{r^3}
=\frac{p}{2\pi\varepsilon_0r^3}.\]

另一同向偶极子的相互作用能为

\[
U_\parallel=-pE_{\rm ax}
=-\frac{p^2}{2\pi\varepsilon_0r^3}.
\]

径向力定义为

\[F_\parallel=-\frac{dU_\parallel}{dr}.\]

求导：

\[
\frac{d}{dr}r^{-3}=-3r^{-4},
\]

所以

\[
F_\parallel
=-\frac{3p^2}{2\pi\varepsilon_0r^4}.
\]

负号表示吸引。\(k\to\infty\) 时 \(r=R\)，故作用力大小

\[
\boxed{|F_\parallel|
=\frac{3\alpha^2E^2}{2\pi\varepsilon_0R^4}}.
\]

\Needspace{3\baselineskip}
\textbf{2. 交变平行外场产生直流项与二倍频项}

外场为

\[E(t)=E_0\cos\omega t,\]

故

\[p(t)=\alpha E_0\cos\omega t.\]

利用

\[\cos^2\omega t=\frac12(1+\cos2\omega t),\]

并在偶极作用中忽略微小 \(\Delta\)，取 \(r\simeq R\)：

\[
\begin{aligned}
F_\parallel(t)
&=-\frac{3\alpha^2E_0^2}{2\pi\varepsilon_0R^4}
\cos^2\omega t\\
&=-\frac{3\alpha^2E_0^2}{4\pi\varepsilon_0R^4}
(1+\cos2\omega t).
\end{aligned}
\]

定义

\[F_0=\frac{3\alpha^2E_0^2}{4\pi\varepsilon_0R^4}.\]

相对运动方程为

\[
\frac m2\ddot\Delta+k\Delta
=-F_0(1+\cos2\omega t).
\]

设稳态特解

\[\Delta(t)=\overline\Delta+A\cos2\omega t.\]

代入方程。常量部分满足

\[k\overline\Delta=-F_0,\]

所以

\[
\boxed{
\overline\Delta_\parallel
=-\frac{3\alpha^2E_0^2}{4\pi\varepsilon_0kR^4}}.
\]

负号表示平均间距缩短。

余弦部分满足

\[
\left[k-\frac m2(2\omega)^2\right]A=-F_0.
\]

因为

\[\frac m2(2\omega)^2=2m\omega^2,\]

故

\[
A=-\frac{F_0}{k-2m\omega^2}.
\]

振幅取绝对值：

\[
\boxed{
|A_\parallel|
=\frac{3\alpha^2E_0^2}
{4\pi\varepsilon_0R^4|k-2m\omega^2|}}.
\]

间距变化的角频率为

\[\boxed{\Omega_{\rm drive}=2\omega}.\]

完整稳态表达式为

\[
\boxed{
\Delta_\parallel(t)
=-\frac{3\alpha^2E_0^2}{4\pi\varepsilon_0kR^4}
-\frac{3\alpha^2E_0^2}
{4\pi\varepsilon_0R^4(k-2m\omega^2)}
\cos2\omega t}.
\]

自由相对振动角频率为

\[\Omega_0=\sqrt{\frac{k}{\mu}}=\sqrt{\frac{2k}{m}}.\]

当 \(2\omega=\Omega_0\)，即

\[k-2m\omega^2=0,\]

无阻尼模型不存在有界稳态，发生共振。

\Needspace{3\baselineskip}
\textbf{3. 垂直外场下的偶极作用}

偶极子在赤道面上产生的电场大小为

\[E_\perp=\frac{p}{4\pi\varepsilon_0r^3},\]

方向与偶极矩相反。因此同向并列偶极子的相互作用能为

\[
U_\perp=+\frac{p^2}{4\pi\varepsilon_0r^3}.
\]

径向力

\[
F_\perp=-\frac{dU_\perp}{dr}
=+\frac{3p^2}{4\pi\varepsilon_0r^4}.
\]

正号表示排斥，大小是同轴吸引力的一半。交变外场下

\[
F_\perp(t)
=+\frac{3\alpha^2E_0^2}{8\pi\varepsilon_0R^4}
(1+\cos2\omega t).
\]

完全类似地得到

\[
\boxed{
\overline\Delta_\perp
=+\frac{3\alpha^2E_0^2}{8\pi\varepsilon_0kR^4}},
\]

\[
\boxed{
|A_\perp|
=\frac{3\alpha^2E_0^2}
{8\pi\varepsilon_0R^4|k-2m\omega^2|}},
\]

\[
\boxed{\Omega_{\rm drive}=2\omega}.
\]

完整表达式为

\[
\boxed{
\Delta_\perp(t)
=\frac{3\alpha^2E_0^2}{8\pi\varepsilon_0kR^4}
+\frac{3\alpha^2E_0^2}
{8\pi\varepsilon_0R^4(k-2m\omega^2)}
\cos2\omega t}.
\]

\Needspace{3\baselineskip}
\textbf{为什么用约化质量}

两球受到大小相等、方向相反的内力。设相对距离 \(r=x_2-x_1\)，则

\[\ddot r=\ddot x_2-\ddot x_1=\frac{F}{m}-\left(-\frac Fm\right)=\frac{2F}{m}.\]

因此相对坐标的惯性系数为
\(m/2\)。若直接对单个小球写方程，弹簧长度改变量是单球位移的两倍，也会得到同一结果。

\Needspace{3\baselineskip}
\textbf{极化反馈为什么被忽略}

题设明确忽略一个原子的电场对另一原子的极化作用，所以始终使用
\(p=\alpha E_{\rm ext}\)。若考虑反馈，则应联立

\[p_1=\alpha(E_{\rm ext}+E_{2\to1}),
\qquad
p_2=\alpha(E_{\rm ext}+E_{1\to2}),\]

偶极矩本身会依赖间距，可能出现更强的非线性耦合。

\Needspace{3\baselineskip}
\textbf{无阻尼题中的``稳态''含义}

本题没有阻尼，通解还含
\(C_1\cos\Omega_0t+C_2\sin\Omega_0t\)。这些自由振动不会自行消失，所以所谓稳态公式是受迫特解，须由适当初值选出，或理解为很小阻尼下充分久后的极限。共振时则不存在有界周期特解。

一般夹角变式还会产生改变分子朝向的力矩。上文
\(F_r\propto1-3\cos^2\beta\)
的结果要求分子轴方向被固定；魔角只消去径向分量，不能据此说全部作用消失。对平行和垂直两个对称方向，转动力矩为零，可采用正文一维模型。

若要检查近似的误差，平行力取 \(r=R+\Delta\) 后
\(F\simeq-F_0(1+\cos2\omega t)(1-4\Delta/R)\)。被忽略的项相对量级为
\(|\Delta|/R\)，在共振附近尤其需要检验。

\Needspace{3\baselineskip}
\textbf{低频、刚性极限与垂直外场的检查}

在 \(\omega\to0\) 的慢变化极限，二倍频特解的系数趋
\(-F_0/k\)，所以平行外场的准静态伸长为

\[\Delta_{\rm qs}(t)=-\frac{F_0}{k}(1+\cos2\omega t).\]

在外场最大时收缩最大，外场为零时回到未极化平衡长度。这个极限须在自由振动已被阻尼压低、或初态按准静态准备的条件下使用；对突然开启的无阻尼模型，自由振动仍在。

垂直外场的力是平行外场的负一半，所以

\[\overline\Delta_\perp=\frac{F_0}{2k},\qquad
A_\perp=\frac{F_0}{2(k-2m\omega^2)}.\]

方向为平均伸长，受迫振幅是平行时的一半。两者在 \(k\to\infty\)
时位移趋零，偶极力本身却不趋零；``刚性键不伸长''不等于``没有相互作用力''。

\Needspace{3\baselineskip}
\textbf{初始静止时，自由振动项不能省略}

设 \(\Omega=\sqrt{2k/m}\)、远离共振。平行外场的方程为

\[\frac m2\ddot\Delta+k\Delta=-F_0(1+\cos2\omega t).\]

若 \(\Delta(0)=\dot\Delta(0)=0\)，完整解是

\[\Delta(t)=-\frac{F_0}{k}-\frac{F_0}{k-2m\omega^2}\cos2\omega t
+\left(\frac{F_0}{k}+\frac{F_0}{k-2m\omega^2}\right)\cos\Omega t.\]

最后一项是为满足初态而加入的自由振动，在无阻尼模型中不会自行消失。正文的``平均伸长''和``受迫振幅''分别对应常数项和二倍频特解；若问任意时刻间距，必须保留完整解。

小位移近似还要求
\(F_0/k\ll R\)、\(F_0/|k-2m\omega^2|\ll R\)，靠近共振时即使电场很弱也可能失效。严格共振的余弦驱动给随时间增长的
\(-F_0t\sin\Omega t/(m\Omega)\)，不能以分母为零解释成一个无限大的稳定振幅。

\Needspace{3\baselineskip}
\textbf{任意固定夹角与转动约束}

若两个同向诱导偶极与键轴夹角为 \(\chi\)，准静态相互作用能为

\[U_{dd}=\frac{p^2}{4\pi\varepsilon_0r^3}(1-3\cos^2\chi).\]

保持方向固定后求径向力，

\[F_r=\frac{3p^2}{4\pi\varepsilon_0r^4}(1-3\cos^2\chi).\]

取 \(\chi=0\) 恢复吸引，取 \(\chi=\pi/2\)
恢复大小为前者一半的排斥。\(\cos^2\chi=1/3\)
时径向项为零，但对角度求导的转矩一般不为零，所以并不代表所有偶极作用消失。

题中只讨论键长振动，应说明键轴方向受限或在所取过程中不变。允许自由转动时，径向与角向运动耦合，不能在运动中始终把
\(\chi\)
当常数。诱导偶极近似还忽略另一原子电场对本原子偶极的自洽反馈，要求
\(\alpha/(4\pi\varepsilon_0R^3)\ll1\)。

\Needspace{3\baselineskip}
\textbf{保留间距反馈后会出现周期刚度}

原解把偶极力中的间距取为 \(R\)，是按外场平方的最低阶近似。若保留
\(r=R+\Delta\)，

\[r^{-4}=R^{-4}(1-4\Delta/R)+O(\Delta^2/R^6).\]

平行场的力变为

\[F_{dd}=-F_0(1+\cos2\omega t)
+\frac{4F_0}{R}(1+\cos2\omega t)\Delta+\cdots.\]

于是方程中刚度是周期变化的：

\[\frac m2\ddot\Delta+
\left[k-\frac{4F_0}{R}(1+\cos2\omega t)\right]\Delta
=-F_0(1+\cos2\omega t).\]

远离共振且 \(\Delta=O(E_0^2)\) 时，反馈项为
\(O(E_0^4)\)，正文最低阶解可忽略它；接近共振时位移增大，则不能以同一小参数次序保证忽略。常值部分使平行场的有效刚度降低，垂直场则增加；因此``力小''与``恢复刚度变化可忽略''也需要分别检查。
\end{document}
